\documentclass[11pt,letterpaper]{article}
\usepackage[T1]{fontenc}
\usepackage[utf8]{inputenc}
\usepackage{lmodern}
\usepackage[margin=1in,headheight=15pt]{geometry}
\usepackage{amsmath,amssymb,amsthm,mathtools,mathrsfs}
\usepackage{microtype,booktabs,longtable,array,enumitem}
\usepackage{xcolor,fancyhdr}
\definecolor{ink}{HTML}{283A35}
\definecolor{accent}{HTML}{49675A}
\usepackage[colorlinks=true,linkcolor=ink,citecolor=accent,urlcolor=accent]{hyperref}
\usepackage{aliascnt}
\usepackage[nameinlink,noabbrev]{cleveref}
\hypersetup{pdftitle={Surquaternions: Algebra, Geometry, Valuation, and Controlled Analysis},pdfauthor={Research exposition prepared for Vladimir Reshetnikov},pdfsubject={Quaternions over the surreal numbers}}
\pagestyle{fancy}\fancyhf{}
\fancyhead[L]{\small\textcolor{ink}{Surquaternions}}
\fancyhead[R]{\small\thepage}
\renewcommand{\headrulewidth}{0.3pt}
\setlength{\parskip}{0.3em}
\setlength{\emergencystretch}{3em}
\setlist{itemsep=0.15em,topsep=0.3em}
\numberwithin{equation}{section}
\newtheorem{theorem}{Theorem}[section]
\newaliascnt{proposition}{theorem}
\newtheorem{proposition}[proposition]{Proposition}
\aliascntresetthe{proposition}
\crefname{proposition}{Proposition}{Propositions}
\newaliascnt{lemma}{theorem}
\newtheorem{lemma}[lemma]{Lemma}
\aliascntresetthe{lemma}
\crefname{lemma}{Lemma}{Lemmas}
\newaliascnt{corollary}{theorem}
\newtheorem{corollary}[corollary]{Corollary}
\aliascntresetthe{corollary}
\crefname{corollary}{Corollary}{Corollaries}
\theoremstyle{definition}
\newaliascnt{definition}{theorem}
\newtheorem{definition}[definition]{Definition}
\aliascntresetthe{definition}
\crefname{definition}{Definition}{Definitions}
\newaliascnt{example}{theorem}
\newtheorem{example}[example]{Example}
\aliascntresetthe{example}
\crefname{example}{Example}{Examples}
\theoremstyle{remark}
\newaliascnt{remark}{theorem}
\newtheorem{remark}[remark]{Remark}
\aliascntresetthe{remark}
\crefname{remark}{Remark}{Remarks}
\newcommand{\No}{\mathbf{No}}
\newcommand{\SQ}{\mathbb H_{\No}}
\newcommand{\HF}{\mathbb H_F}
\newcommand{\HR}{\mathbb H_{\mathbb R}}
\newcommand{\R}{\mathbb R}
\newcommand{\C}{\mathbb C}
\newcommand{\Q}{\mathbb Q}
\newcommand{\Z}{\mathbb Z}
\newcommand{\N}{\mathbb N}
\newcommand{\OO}{\mathcal O}
\newcommand{\mm}{\mathfrak m}
\newcommand{\Oz}{\mathbf{Oz}}
\newcommand{\eps}{\varepsilon}
\newcommand{\ii}{\mathbf i}
\newcommand{\jj}{\mathbf j}
\newcommand{\kk}{\mathbf k}
\newcommand{\Rea}{\operatorname{Re}}
\newcommand{\Ima}{\operatorname{Im}}
\newcommand{\st}{\operatorname{st}}
\newcommand{\supp}{\operatorname{supp}}
\newcommand{\lc}{\operatorname{lc}}
\newcommand{\tr}{\operatorname{tr}}
\newcommand{\Aut}{\operatorname{Aut}}
\newcommand{\GL}{\operatorname{GL}}
\newcommand{\SO}{\operatorname{SO}}
\newcommand{\ad}{\operatorname{ad}}
\newcommand{\diag}{\operatorname{diag}}
\newcommand{\ord}{\operatorname{ord}}
\newcommand{\Exp}{\operatorname{Exp}}
\newcommand{\Log}{\operatorname{Log}}
\newcommand{\BCH}{\operatorname{BCH}}
\newcommand{\Span}{\operatorname{span}}
\newcommand{\abs}[1]{\left|#1\right|}
\newcommand{\norm}[1]{\left\|#1\right\|}
\newcommand{\angles}[1]{\left\langle#1\right\rangle}
\newcommand{\DD}{\mathcal D}
\newcommand{\id}{\operatorname{id}}
\newcolumntype{L}[1]{>{\raggedright\arraybackslash}p{#1}}
\begin{document}
\begin{titlepage}
\centering
\vspace*{0.5cm}
{\large\scshape Algebra at every surreal scale\par}
\vspace{0.7cm}
{\Huge\bfseries\textcolor{ink}{Surquaternions}\par}
\vspace{0.35cm}
{\Large Algebra, Geometry, Valuation,\\[0.2em]and Controlled Analysis\par}
\vspace{0.7cm}
{\large September 21, 2026\par}
\vspace{0.6cm}
\begin{minipage}{0.95\textwidth}
\begin{abstract}
A surquaternion is an element of
$\SQ=\No\oplus\No\ii\oplus\No\jj\oplus\No\kk$ with Hamilton's multiplication.
The definition is elementary; a usable theory must also explain which parts of quaternionic
geometry survive non-Archimedean scales, what an infinite series means, and which notions of
analysis are compatible with the surreal field.

We develop the division algebra, its algebraically closed complex slices, conjugacy classes,
rotation groups, one-sided polynomial equations, and quaternionic spectral theory over real
closed fields. We identify surquaternions with set-supported Hahn series having ordinary
quaternion coefficients, prove the valuation and standard-part formulas, and separate that
summation theory from the exceptionally fine topology of the full surreal class. Further
results include a complete polynomial zero-class algorithm, an infinitesimal spectral-projector
estimate with a sharp two-scale example, a noncommutative exponential--logarithm correspondence
near the identity, and a coefficientwise slice and Fueter calculus.

A central compatibility theorem states that the Berarducci--Mantova derivation, extended
componentwise, admits no nonzero solution of $Dy=\ii y$ in $\SQ$. Consequently no global
quaternionic exponential can satisfy the usual commuting chain rule for that derivation at
$\omega\ii$. This does not prevent a global exponential: the Ehrlich--Kaplan trigonometric
normalization supplies one, but it is a different analytic structure. We specify this distinction
rather than suppressing it.
\end{abstract}
\end{minipage}
\vfill
\begin{minipage}{0.95\textwidth}\small
\textbf{Status.} A self-contained mathematical development, with cited foundational inputs and
proofs of the stated extensions. This is not a refereed paper, a priority claim, or a claim to
solve a named published open problem. The accompanying symbolic checks verify finite identities
and examples, not all theorems. No Lean verification of this article is claimed.

\textbf{Repository baseline.} The source repository is Vladimir Reshetnikov's \texttt{Surreal},
read at commit \texttt{39f2be6667ade51bca2b45daa47e289d69c09764}.
The article is a new standalone companion, not a modification of that repository.
\end{minipage}
\end{titlepage}
\tableofcontents
\clearpage

\section{Scope, foundations, and the choice of analytic structure}\label{sec:scope}

\subsection{What is being generalized?}
Hamilton's quaternions do not require real completeness for their multiplication. They require
an ordered field to make a sum of four squares positive, and a real closed field to obtain the
square roots and complex factorizations used in much of their geometry. The surreal numbers
supply both, while also supplying infinitely large and infinitesimal coefficients.

The basic object is therefore
\[
 \SQ=\{a+b\ii+c\jj+d\kk:a,b,c,d\in\No\},\qquad
 \ii^2=\jj^2=\kk^2=\ii\jj\kk=-1.
\]
We use \emph{surquaternion} descriptively for this object, without asserting that the word or
construction is new. Throughout, coefficients commute with $\ii,\jj,\kk$; this is not an
algebra of games with a new multiplication, nor an attempt to order the quaternion algebra.

Three levels organize the development. Finite algebra is valid over a real closed field $F$.
Hahn algebra adds well-ordered, set-sized supports and exact infinite summation. Analytic
constructions additionally specify domains, support conditions, and the differentiation being
used. A theorem at one level is not silently promoted to another.

\subsection{Relation to the supplied repository}
The repository's existing documents distinguish the ordered-field modulus, Hahn valuation,
standard part, and coefficientwise analytic constructions. Its trigonometry report emphasizes
that directions at arbitrary scale have finite angles, while also recording global
Ehrlich--Kaplan trigonometry. We retain these distinctions
\cite{Repo,RepoDocs,RepoTrig}.

The repository's Lean implementation already develops generic finite algebra and Hahn support
infrastructure; its README explicitly says that it does not yet construct the surreal field or
the normal-form bridge. We neither identify those prerequisites with a formalization of this
article nor infer a new theorem from a successful build reported elsewhere. The proposed
formalization layers in \cref{sec:computation} concern additional work.

Much of the algebra below is classical quaternion algebra over a real closed field, written out
to make its hypotheses visible. The slice framework and Fueter transformation have classical
antecedents \cite{Voight,GP}. The genuinely surreal issues are the class/set boundary, exact
support control, multiscale estimates, and incompatibility between otherwise natural analytic
structures. Even a useful derived theorem in these directions is not evidence of bibliographic
priority; the literature review here is targeted, not exhaustive.

\subsection{Foundational inputs and set-sized workspaces}
We use the established facts that $\No$ is a real closed ordered class field, that it has Conway
normal forms, and that Gonshor's exponential is an ordered exponential-field operation
\cite{Conway,Gonshor,vdDE}. Class language can be interpreted in a class theory such as NBG;
alternatively each set of parameters can be placed in the following set-sized workspace.
For the full surreal class, topological statements use the displayed ball basis, not a
postulated set of all open subclasses. Algebraic group points are finite coordinate tuples.
Statements about arbitrary class automorphisms or derivations concern individual class maps;
we do not postulate a class whose elements are all class functions. Set-field theorems about
finite matrices and polynomials pass to their surreal instances through the workspace argument.

Write
\begin{equation}\label{eq:monomial}
 t^\gamma:=\omega^{-\gamma}.
\end{equation}
Here the right side is the \emph{Conway monomial map}. It must not be replaced, for arbitrary
surreal $\gamma$, by a presumed identity with
$\exp(-\gamma\log\omega)$ involving Gonshor's exponential. The two notational uses of powers
must be kept distinct.

For a set-sized ordered subgroup $\Gamma\subset(\No,+)$ let
\[
 K_\Gamma=\R((t^\Gamma)),\qquad D_\Gamma=\mathbb H_{K_\Gamma}.
\]
A Hahn series has a well-ordered support in $\Gamma$, with our increasing-exponent convention.
If $\Gamma$ is divisible, the standard real-closed Hahn-field theorem makes $K_\Gamma$ real
closed. This theorem, and the surreal normal-form theorem that embeds these series in $\No$,
are foundational inputs, not consequences of quaternion multiplication.

\begin{theorem}[Set localization]\label{thm:localization}
Every set $A$ of surquaternions is contained in some $D_\Gamma$ with $\Gamma$ a divisible,
set-sized ordered subgroup of $(\No,+)$. Its elements are represented by their actual Conway
normal forms, not merely by an abstract field embedding.
\end{theorem}
\begin{proof}
Take the union of the supports of all four coordinates of all elements of $A$. This is a set
of surreal exponents. Its rational linear span is a set-sized divisible subgroup $\Gamma$.
All the coordinate series belong to $K_\Gamma$, so all the original quaternions belong to
$D_\Gamma$. The normal-form addition and multiplication identify the operations there with
those in $\SQ$.
\end{proof}

This localization does \emph{not} assert that $K_\Gamma$ is stable under exponential,
a chosen surreal derivation, or every analytic operation later considered. Specified images
can be included in a larger workspace. Nor does localization identify the workspace's intrinsic
topology with its subspace topology in the full surreal class; these are dramatically different.

For fixed finite matrix sizes and polynomial degrees, algebraic statements about $\HF$ can be
translated into ordered-field statements about four coordinates per variable. First-order
real-closed-field transfer is therefore available for such statements. It is not a license to
transfer compactness, infinite summation, arbitrary function spaces, or the proper-class
assertions about $\No$. Most proofs below avoid even this shortcut.

\section{The division algebra and its norm}\label{sec:algebra}

Let $F$ initially be any ordered field. We write $q=(a,u)$, where $a\in F$ and $u\in F^3$.
The multiplication law is
\begin{equation}\label{eq:product}
 (a,u)(b,v)=(ab-u\cdot v,\;av+bu+u\times v).
\end{equation}
The dot and cross products are the polynomial formulas with the usual orientation
$(\ii,\jj,\kk)$. In particular no metric completeness is hidden in this notation.

Define
\[
 \bar q=(a,-u),\qquad T(q)=q+\bar q=2a,\qquad
 N(q)=q\bar q=a^2+u\cdot u.
\]

\begin{theorem}[Hamilton algebra over an ordered field]\label{thm:division}
The algebra $\HF$ is an associative division algebra with center $F$. Conjugation is an
anti-involution, and
\begin{equation}\label{eq:inverse}
 N(pq)=N(p)N(q),\qquad q^{-1}=\frac{\bar q}{N(q)}\quad(q\ne0).
\end{equation}
Every element satisfies
\begin{equation}\label{eq:quadratic}
 q^2-T(q)q+N(q)=0.
\end{equation}
\end{theorem}
\begin{proof}
Associativity can be checked from \eqref{eq:product}, or from the faithful matrix
representation below. The same formula gives $\overline{pq}=\bar q\bar p$ and
$q\bar q=\bar q q=N(q)$. If $q\ne0$, at least one coordinate square is positive, so $N(q)>0$.
The inverse formula follows. Consequently
$N(pq)=pq\bar q\bar p=N(q)p\bar p=N(p)N(q)$.
An element commuting with both $\ii$ and $\jj$ has zero imaginary coordinates; hence the
center is exactly $F$. Expanding $q(q+\bar q)$ proves \eqref{eq:quadratic}.
\end{proof}

In particular $\SQ$ has no zero divisors. It cannot carry an order compatible with its ring
operations: such an order would make every nonzero square positive, whereas $\ii^2=-1$.
Its ordered geometry comes from $F$, not from an order on $\HF$ itself.

Suppose from now on that $F$ is real closed, unless a weaker hypothesis is stated. Define the
\emph{$F$-valued modulus} and scalar inner product by
\[
 |q|=\sqrt{N(q)},\qquad (p,q)=\Rea(p\bar q).
\]
This is not an ordinary real norm when $F$ is non-Archimedean.

\begin{proposition}[Euclidean identities without completeness]\label{prop:norm}
For $p,q\in\HF$,
\[
 |pq|=|p||q|,\qquad |p+q|\le |p|+|q|,
 \qquad \max_j|q_j|\le |q|\le2\max_j|q_j|.
\]
For nonzero $p,q$, equality in the triangle inequality holds exactly when $q=\lambda p$
for some $\lambda\in F_{>0}$.
\end{proposition}
\begin{proof}
The first and third assertions follow from multiplicativity of $N$ and the four coordinate
squares. For Cauchy--Schwarz, the polynomial
$N(p-sq)$ is nonnegative for every $s\in F$. Substituting
$s=(p,q)/N(q)$ gives $(p,q)^2\le N(p)N(q)$.
Therefore $N(p+q)\le(|p|+|q|)^2$, proving the triangle inequality.
Equality forces $(p,q)=|p||q|$ and equality in Cauchy--Schwarz; the same substitution gives
$p=((p,q)/N(q))q$ with positive scalar coefficient. The converse is immediate.
\end{proof}

\subsection{Complex coordinates and splitting}
Put $C_F=F[\ii]$. Write $q=z+w\jj$ with $z,w\in C_F$. Since
$\jj z=\bar z\jj$, multiplication becomes
\[
 (z+w\jj)(z'+w'\jj)
 =(zz'-w\bar w')+(zw'+w\bar z')\jj.
\]
Thus
\begin{equation}\label{eq:rho}
 \rho(z+w\jj)=
 \begin{pmatrix}z&w\\-\bar w&\bar z\end{pmatrix}
\end{equation}
is an injective algebra homomorphism, with
\[
 \rho(\bar q)=\rho(q)^*,\qquad
 \det\rho(q)=N(q),\qquad \tr\rho(q)=T(q).
\]
Associativity follows as well. Scalar extension gives
\[
 \HF\otimes_F C_F\simeq M_2(C_F).
\]
The image of $\HF$ in \eqref{eq:rho} is a special four-dimensional real form of the matrix
algebra, not the entire matrix algebra. In particular, complex scalar extension splits the
algebra and introduces zero divisors; it does not exhibit zero divisors in $\HF$.

More generally, a quaternion symbol algebra $(a,b)_F$ over a real closed field is split unless
both $a$ and $b$ are negative. Indeed a positive parameter is a square and produces a
nontrivial idempotent, whereas two negative parameters can be rescaled to $(-1,-1)_F$.
Thus Hamilton's algebra is the quaternion division-algebra case. This modest classification
statement is sufficient here; no classification of arbitrary class-sized division algebras is
being invoked.

\section{Complex slices, conjugacy, and scalar equations}\label{sec:slices}

Define
\[
 \mathbb S(F)=\{I\in\HF:\Rea I=0,\ N(I)=1\}.
\]
Equation \eqref{eq:quadratic} shows that this is exactly the set of solutions of $I^2=-1$.
For $I\in\mathbb S(F)$, the subalgebra $F_I=F+FI$ is isomorphic to $F[\ii]$ and is therefore
algebraically closed when $F$ is real closed.

Every nonreal quaternion has a unique expression
\begin{equation}\label{eq:slice-decomp}
 q=a+Ir,\qquad a=\Rea q,\quad r=|\Ima q|>0,\quad
 I=\frac{\Ima q}{|\Ima q|}.
\end{equation}
Two such slice fields coincide precisely when their units differ by sign, and otherwise
intersect in $F$. The whole quaternion algebra is their union, but it is not their disjoint
union and multiplication need not stay in either of two different slices.

\begin{theorem}[Centralizers and conjugacy]\label{thm:conjugacy}
For nonreal $q=a+Ir$, its centralizer is $F_I$. Two quaternions are conjugate by an invertible
quaternion if and only if they have the same trace and norm. In particular, the conjugacy
class of $a+Ir$ is
\[
 [a,r]=a+r\mathbb S(F).
\]
Every commuting pair lies in a common complex slice.
\end{theorem}
\begin{proof}
The commutator of $(a,u)$ and $(b,v)$ is $(0,2u\times v)$. For $u\ne0$, its vanishing is
equivalent to $v$ being an $F$-multiple of $u$, proving the centralizer assertion.
Conjugation preserves scalar part and norm. Conversely, let $I,J\in\mathbb S(F)$.
If $J\ne-I$, then $g=1-JI\ne0$ satisfies
$gI=I+J=Jg$. If $J=-I$, choose a unit imaginary $K$ perpendicular to $I$; then
$KI=-IK=JK$. These elements conjugate $I$ to $J$. The last assertion follows from the
centralizer description, with any slice available for two scalar elements.
\end{proof}

A nonreal surquaternion thus has a proper-class sphere of conjugates, even though its minimal
polynomial over the center has degree two. Polynomial degree bounds must count conjugacy
classes or multiplicities, not individual points on every such sphere.

\begin{proposition}[Square roots]\label{prop:sqrt}
For $q=a+v$ not on the nonpositive scalar axis, let
\[
 s=\sqrt{\frac{|q|+a}{2}}.
\]
Then $s>0$ and the two square roots of $q$ are
\begin{equation}\label{eq:sqrt}
 \pm\left(s+\frac{v}{2s}\right).
\end{equation}
A negative scalar $-\rho$, $\rho>0$, has the roots
$\sqrt\rho\,I$ for all $I\in\mathbb S(F)$. A positive scalar has just its two scalar
square roots, and zero has only zero.
\end{proposition}
\begin{proof}
The square of the displayed expression has imaginary part $v$ and real part
$s^2-|v|^2/(4s^2)=a$, using $|v|^2=|q|^2-a^2$.
Any square root commutes with its square. For nonreal $q$ it therefore belongs to $F_I$,
where a quadratic has only two roots. For scalar $q$, writing a prospective root $b+u$
gives imaginary part $2bu$. If $u=0$ one obtains the scalar roots; if $b=0$ its square is
$-|u|^2$, giving exactly the negative-scalar sphere.
\end{proof}

For example, $2+\ii$ squares to $3+4\ii$, while
$\omega^{-1}\ii+\omega^{-2}\jj$ has a square-root scale determined by the positive square
root of its norm. No choice of an ``infinite angle'' is required to solve these equations.

\section{Quaternion-valued Hahn series}\label{sec:hahn}

\subsection{Normal form and multiplication}
Combining four Conway normal forms gives a unique expansion
\begin{equation}\label{eq:qnormal}
 q=\sum_{\gamma\in S}h_\gamma t^\gamma,
 \qquad h_\gamma\in\HR\setminus\{0\},
\end{equation}
where $S\subset\No$ is a well-ordered set. Uniqueness follows coordinatewise; a finite union
of well-ordered supports is well ordered. On a workspace this identifies
\begin{equation}\label{eq:hahnbridge}
 D_\Gamma\simeq\HR((t^\Gamma)).
\end{equation}
The monomials are central. Consequently this is a Hahn series algebra with noncommutative
coefficients, not a series construction with noncommuting exponents.

A family $(q_\lambda)$ is \emph{strongly summable} if the union of its supports is well ordered
and each exponent occurs in only finitely many members. Its sum is defined coefficientwise.
This is the usual surreal/Hahn summability notion \cite{BM}; coefficients at a fixed exponent
are added only finitely many times. In particular ordinary convergence of infinitely many
real coefficients is not a replacement for this condition.

For products,
\[
 \left(\sum_\alpha h_\alpha t^\alpha\right)
 \left(\sum_\beta k_\beta t^\beta\right)
 =\sum_\gamma\left(\sum_{\alpha+\beta=\gamma}h_\alpha k_\beta\right)t^\gamma.
\]
Well-ordered-support arithmetic makes the inner sums finite. The order of $h_\alpha$ and
$k_\beta$ is essential, but finite coefficient rearrangement respecting this order proves
associativity and distributivity.

\begin{theorem}[Valuation and residue]\label{thm:valuation}
For $q\ne0$, define $w(q)=\min\supp(q)$ and $\lc(q)=h_{w(q)}$, and put $w(0)=+\infty$.
Then
\begin{align}
 w(pq)&=w(p)+w(q),&\lc(pq)&=\lc(p)\lc(q),\label{eq:valuation-product}\\
 w(p+q)&\ge\min\{w(p),w(q)\},&w(\bar q)&=w(q),\\
 w(N(q))&=2w(q),&w(|q|)&=w(q).\label{eq:normval}
\end{align}
The valuation ring and its maximal two-sided ideal are
\[
 \OO_H=\{q:w(q)\ge0\},\qquad \mm_H=\{q:w(q)>0\}.
\]
They consist respectively of quaternions with finite coordinates and with infinitesimal
coordinates. Coefficient-zero extraction induces
\begin{equation}\label{eq:residue}
 \st:\OO_H\longrightarrow\HR,\qquad
 \ker\st=\mm_H,\qquad \OO_H/\mm_H\simeq\HR.
\end{equation}
\end{theorem}
\begin{proof}
At the smallest possible exponent in a product, the coefficient is the product of the two
leading coefficients. That product is nonzero because $\HR$ is a division algebra. This
proves \eqref{eq:valuation-product}; the ultrametric inequality is coefficientwise.
The leading coefficient of $N(q)=q\bar q$ is $N(\lc(q))>0$, giving the first equality in
\eqref{eq:normval}; taking its positive square root gives the second.
A negative leading exponent is infinitely large, a positive one infinitesimal, and exponent
zero has a real quaternion coefficient followed by infinitesimal terms. On nonnegative
supports only two zero exponents can contribute to the zero coefficient of a product.
Thus $\st$ is a ring homomorphism with the indicated kernel and image.
\end{proof}

In particular $\OO_H^\times$ consists exactly of elements with nonzero standard part.
Valuation is invariant under inner conjugation, while its leading quaternion is conjugated
by the leading quaternion of the conjugating element. The associated graded algebra in a
workspace is the group-graded algebra with one copy of $\HR$ at each $\gamma\in\Gamma$;
it is naturally written $\HR[\Gamma]$ with central monomials.

\subsection{Support-controlled inverse and leading-term decomposition}
We use the Neumann support lemma in its strong form: if $S$ is a well-ordered set of strictly
positive elements of an ordered abelian group, the monoid of finite sums from $S$ is well
ordered, and any fixed exponent has only finitely many representations by words from $S$,
across all word lengths. This is the combinatorial input behind infinitesimal power-series
evaluation; its use in surreal summation is discussed in \cite{BM}, and the repository
implements a corresponding Hahn support theorem \cite{Repo}.

\begin{proposition}[Exact geometric inverse]\label{prop:neumann}
Every nonzero $q$ has a unique decomposition
\begin{equation}\label{eq:leadingdecomp}
 q=t^\gamma h(1+\eps),\qquad
 \gamma=w(q),\ h=\lc(q)\in\HR^\times,\ \eps\in\mm_H.
\end{equation}
Its inverse is the strong Hahn sum
\begin{equation}\label{eq:neumann}
 q^{-1}=\left(\sum_{n\ge0}(-\eps)^n\right)h^{-1}t^{-\gamma}.
\end{equation}
For nonzero $\eps\in\mm_H$ and $M\ge0$,
\begin{equation}\label{eq:neumannerror}
 w\left((1+\eps)^{-1}-\sum_{n=0}^M(-\eps)^n\right)
 =(M+1)w(\eps).
\end{equation}
\end{proposition}
\begin{proof}
Set $\eps=h^{-1}t^{-\gamma}q-1$. The leading term has been removed, so all remaining
exponents are positive. The support lemma makes the geometric family strongly summable.
Multiplication by $1+\eps$ telescopes coefficientwise to $1$, on either side. Its finite
remainder is exactly $(-\eps)^{M+1}(1+\eps)^{-1}$; the final factor has valuation zero.
The valuation identities prove \eqref{eq:neumannerror}.
\end{proof}

The order in \eqref{eq:neumann} is not cosmetic: in general $h$ and $\eps$ do not commute.
The multiplicative structure is a central monomial factor together with a semidirect product
of principal units and ordinary nonzero quaternions. The principal units are already
nonabelian; \cref{sec:explog} describes their Lie structure.

\begin{example}[Two scales in one inverse]
Let $t=\omega^{-1}$ and $q=1+t\ii+t^2\jj$. Then
\[
 q^{-1}=\frac{1-t\ii-t^2\jj}{1+t^2+t^4}.
\]
Its formal geometric expansion agrees exactly with \eqref{eq:neumann}. The terms of the
inverse are not obtained by independently inverting the four coordinate series. The absence
of a $\kk$ term here is the identity $\ii\jj+\jj\ii=0$, not an approximation.
\end{example}

\section{Why ordinary topological analysis is insufficient}\label{sec:topology}

The \emph{fine topology} on $\SQ$ has basic balls
$B(q,r)=\{p:|p-q|<r\}$ for every positive surreal $r$. Its restriction to an ordinary
Euclidean-looking set is much finer than the topology suggested by a drawing.

\begin{theorem}[Set-discreteness and the failure of sequential completion]\label{thm:discrete}
Every set $A\subset\SQ$ is uniformly separated and closed in the fine topology. Every
set-indexed fine-Cauchy net is eventually constant, and every convergent set-indexed net is
eventually equal to its limit. Every continuous map from an ordinary connected real interval
to $\SQ$, with this topology, is constant.
\end{theorem}
\begin{proof}
The positive distances between distinct points of $A$ form a set of positive surreals. The
surreal cut property supplies a positive $r$ with $2r$ smaller than all of them. Thus every ball of
radius $r$ contains at most one point of $A$. For a point $q\notin A$, the same argument
applied to the set $\{|q-a|:a\in A\}$ supplies a ball at $q$ disjoint from $A$. This proves
closedness and discreteness, including the empty and singleton cases trivially.

A set-indexed net has a set as its image. The Cauchy condition applied to a uniform separation
radius forces all sufficiently late terms to coincide. If a net converges, apply the
separation argument to its image together with the proposed limit; convergence then forces
eventual equality to that limit. Finally a continuous interval map has a set-sized image,
which is discrete in the induced topology. A connected subset of a discrete space has at
most one point.
\end{proof}

This theorem is about the \emph{full} surreal class and its fine topology. It is false as a
blanket assertion about every useful set-sized Hahn field with its own topology.
For example, in $\R((t^\Q))$ with its intrinsic valuation topology, $t^n\to0$, because the
integers are cofinal in $\Q$. The geometric partial sums converge there to $(1-t)^{-1}$.
Inside $\No$, however, every $t^n$ is larger than $t^\omega$; neither assertion of convergence
survives the finer neighborhood test. The subspace topology on the entire set-sized
workspace inherited from $\No$ is discrete, whereas its intrinsic topology is not.

Even intrinsically, $n\delta$ need not be cofinal in a higher-rank value group when
$\delta>0$. Thus a remainder estimate $w(R_n)\ge n\delta$ is not automatically a statement
that $R_n\to0$ in that group's valuation topology. Strong Hahn summation avoids both
ambiguities by specifying its coefficients directly.

\begin{example}[Zero fine derivative does not imply constancy]\label{ex:zeroderivative}
On $\SQ$, let $f(q)=0$ when all coordinates of $q$ are finite and $f(q)=1$ otherwise.
The finite-coordinate part is open: a ball of radius $1$ around one of its points stays
finite. Its complement is open: at an infinite $q$, the ball of radius $|q|/2$ stays infinite
by the reverse triangle inequality. Hence $f$ is locally constant and has zero differential
at every point, but it is not globally constant.
\end{example}

The example is not a paradox. A mean-value argument would require a connectedness or
integration theorem that this topology does not provide. Similarly, the coefficientwise
Cauchy formula proved later is not a limit of fine-topological Riemann sums along a
nonconstant fine-continuous ordinary path.

For the rest of the article, an infinite expression is therefore marked by its construction:
strong Hahn summation, ordinary real or complex analysis on coefficient functions, or a
specified surreal operation. The word ``convergent'' is not used as a substitute for these
choices.

\section{Rotations, automorphisms, and rational geometry}\label{sec:geometry}

Let $V=\Ima\HF\simeq F^3$ and
\[
 S^3(F)=\{u\in\HF:N(u)=1\}.
\]
Every $q\ne0$ has the unique polar factorization $q=|q|u$ with $u\in S^3(F)$. Because the
positive scalar factor is central,
\begin{equation}\label{eq:polargroup}
 \HF^\times\simeq F_{>0}\times S^3(F)
\end{equation}
as groups. This is algebraic polar decomposition; it does not assert topological compactness
of $S^3(\No)$.

\subsection{The rotation double cover as an algebraic statement}
For a unit quaternion $u=a+v$, conjugation on $V$ is
\begin{equation}\label{eq:rodrigues}
 R_u(x)=uxu^{-1}
 =(a^2-|v|^2)x+2(v\cdot x)v+2a(v\times x).
\end{equation}
The same right side equals $ux\bar u$ without the unit hypothesis. This is an exact polynomial
identity at arbitrary surreal scales.

\begin{theorem}[Rotations and inner automorphisms]\label{thm:rotations}
Over a real closed field $F$, the map $u\mapsto R_u$ gives an exact sequence of groups
\begin{equation}\label{eq:spin3}
 1\longrightarrow\{1,-1\}\longrightarrow S^3(F)
 \longrightarrow\SO(3,F)\longrightarrow1.
\end{equation}
Moreover
\[
 \Aut_F(\HF)\simeq\SO(3,F)\simeq\HF^\times/F^\times.
\]
In particular every center-fixing automorphism of the surquaternions is inner.
\end{theorem}
\begin{proof}
Conjugation preserves $V$, its norm, and its cross product, so its matrix is in $\SO(3,F)$.
Its kernel commutes with every imaginary quaternion and is therefore scalar; unit scalars are
$\pm1$.

For surjectivity, let $A\in\SO(3,F)$. Orthogonality and determinant $1$ give
$\det(A-I)=-\det(A-I)$, so $A$ has a nonzero fixed vector. Normalize it to $I\in V$.
On $I^\perp$ the transformation is a two-dimensional rotation with entries $c,s\in F$,
$c^2+s^2=1$, in an appropriately oriented orthonormal basis. If $c\ne-1$, choose
$a=\sqrt{(1+c)/2}$ and $b=s/(2a)$. Then $a^2+b^2=1$, and \eqref{eq:rodrigues} shows that
$a+bI$ induces $A$. If $c=-1$, use $I$ itself.

An $F$-algebra automorphism preserves the coefficients of the quadratic minimal polynomial
of every nonreal element. It therefore preserves scalar part and norm, and preserves the
imaginary dot and cross products by \eqref{eq:product}. Its action on $V$ lies in $\SO(3,F)$.
Conversely each such orthogonal transformation preserves \eqref{eq:product}, hence extends
by fixing $F$ to an algebra automorphism. Surjectivity just proved makes it inner.
\end{proof}

Equation \eqref{eq:spin3} is a two-to-one map on algebraic group points. It is not, without a
separate topology and definition, a covering-space theorem for the fine surreal topology.
An arbitrary ring automorphism may also act on the center. Restriction to $F$ and the
coefficientwise extension of any field automorphism show that, for set-sized $F$, the full
automorphism group is a semidirect product of the inner automorphisms with $\Aut(F)$.
For $\No$, this statement is read for each individual class automorphism, without forming
a collection of all class functions as an additional mathematical object.

The same approach yields the four-dimensional rotation action
\begin{equation}\label{eq:spin4}
 S^3(F)\times S^3(F)\longrightarrow\SO(4,F),\qquad
 (u,v):q\longmapsto uqv^{-1},
\end{equation}
with kernel $\{(1,1),(-1,-1)\}$. To prove surjectivity, take an orientation-preserving
isometry $A$ of $\HF$ and put $z=A(1)$. Left multiplication by $z^{-1}$ has determinant $1$
and makes the composite fix $1$. Its restriction to $V$ belongs to $\SO(3,F)$, so equals
conjugation by a unit $u$. Thus $A(q)=zuqu^{-1}$, as required. Here the determinant of left
multiplication by a quaternion is $N(q)^2$, obtainable directly from its four-coordinate
matrix.

\subsection{Standard part of rotations and a rational chart}
Every coordinate of a unit surquaternion has absolute value at most $1$. Standard part thus
gives a split exact sequence
\begin{equation}\label{eq:unitsreduction}
 1\longrightarrow K\longrightarrow S^3(\No)
 \xrightarrow{\st} S^3(\R)\longrightarrow1,
 \qquad K=S^3(\No)\cap(1+\mm_H).
\end{equation}
The section is the inclusion of constant quaternions. Its kernel describes infinitesimal
rotations, and is nonabelian. In \cref{thm:principal-log} it will be identified, as a set with
an explicitly transported group law, with the infinitesimal imaginary quaternions.

There is also a purely rational parametrization, requiring no trigonometric functions:
\begin{equation}\label{eq:cayley}
 \operatorname{Cay}(x)=(1+x)(1-x)^{-1},\qquad x\in V.
\end{equation}
Since $N(1-x)=1+N(x)>0$, it is always defined. Conjugation shows that its norm is $1$, and
it never equals $-1$. Conversely, for $u\in S^3(F)\setminus\{-1\}$,
\[
 x=(u-1)(u+1)^{-1}
\]
is imaginary and inverts \eqref{eq:cayley}. All factors here are rational functions of one
quaternion and therefore commute with one another where that is used in the verification.
This gives coordinates on the unit sphere at every scale, including neighborhoods whose
radii are infinitesimal.

\subsection{Quaternionic fractional linear transformations}
A useful further geometry is the right projective line $\mathbb P^1(\HF)$. A matrix
$M=\left(\begin{smallmatrix}a&b\\c&d\end{smallmatrix}\right)\in\GL_2(\HF)$ acts by
\begin{equation}\label{eq:mobius}
 f(q)=(aq+b)(cq+d)^{-1}
\end{equation}
where the denominator is nonzero, with the usual projective interpretation at its pole and at
infinity. Right lines are essential: the column $(q,1)^T$ is multiplied on the left by $M$,
then normalized on the right by $(cq+d)^{-1}$.

For finite affine arguments in its domain, a direct subtraction gives
\begin{equation}\label{eq:mobiusdiff}
 (f(q)-f(p))(cp+d)=(a-f(q)c)(q-p).
\end{equation}
Its algebraic differential is therefore
\[
 df_q(h)=(a-f(q)c)h(cq+d)^{-1}.
\]
The first factor is nonzero: otherwise the row $(a,b)$ would equal $f(q)(c,d)$, contradicting
invertibility of $M$. Consequently
\[
 |df_q(h)|=\frac{|a-f(q)c|}{|cq+d|}|h|.
\]
These transformations are conformal in this exact algebraic differential sense over $F$.
No real compactification, manifold completeness, or contour integration is needed for the
claim. Formula \eqref{eq:mobiusdiff} also illustrates why moving a denominator to the wrong
side changes the transformation.

\section{One-sided polynomials: roots, spheres, and multiplicity}\label{sec:polynomials}

Use a \emph{central} indeterminate $X$ and write
\[
 P(X)=\sum_{k=0}^nX^k a_k\in\HF[X],\qquad
 P(q)=\sum_{k=0}^nq^k a_k.
\]
Coefficients are on the right in evaluation. Denote the multiplication in this polynomial
ring by $\star$ when evaluation is nearby. It is not pointwise multiplication of the
associated quaternion-valued functions.

\begin{lemma}[Twisted product and division]\label{lem:twisted}
If $P(q)\ne0$, then
\begin{equation}\label{eq:twisted}
 (P\star Q)(q)=P(q)\,Q\bigl(P(q)^{-1}qP(q)\bigr).
\end{equation}
If $P(q)=0$, then $(P\star Q)(q)=0$. For any $q$ there is a unique polynomial $Q$ with
\begin{equation}\label{eq:leftdivision}
 P=(X-q)\star Q+P(q).
\end{equation}
Thus a zero $q$ gives a \emph{left} linear factor when coefficients are evaluated on the right.
\end{lemma}
\begin{proof}
Writing $Q=\sum_jX^jb_j$, the product evaluates to
$\sum_jq^jP(q)b_j$. If $P(q)\ne0$, use
$q^jP(q)=P(q)(P(q)^{-1}qP(q))^j$; if $P(q)=0$ each term vanishes.
For division, solve coefficients downward: $b_{n-1}=a_n$ and
$b_{k-1}=a_k+qb_k$. The final remainder is
$a_0+qb_0=\sum_kq^ka_k$. This also proves uniqueness.
\end{proof}

Let $P^c=\sum X^k\bar a_k$ and define the \emph{normal polynomial}
\begin{equation}\label{eq:normalpoly}
 N_P=P\star P^c=P^c\star P\in F[X].
\end{equation}
To see centrality explicitly, write
$P=p_0+p_1\ii+p_2\jj+p_3\kk$ with $p_j\in F[X]$. Then
\[
 N_P=p_0^2+p_1^2+p_2^2+p_3^2.
\]
For degree $n$ its degree is $2n$, with leading coefficient $N(a_n)>0$.
Also $N_{P\star Q}=N_PN_Q$. At a real scalar $x$, $N_P(x)=N(P(x))$; at a nonreal
quaternion, one must not replace the polynomial normal by that pointwise formula.

For a nonreal conjugacy class $[a,r]$, set
\[
 \Delta_{a,r}(X)=X^2-2aX+a^2+r^2,\qquad r>0.
\]
This is irreducible over $F$, and vanishes at every point of that class. Divide centrally:
\begin{equation}\label{eq:sphereremainder}
 P=\Delta_{a,r}Q+XA+B,\qquad A,B\in\HF.
\end{equation}

\begin{theorem}[Complete zero-class criterion]\label{thm:polynomialzeros}
Let $F$ be real closed and let $P$ be nonzero. For each $a\in F$, $r>0$:
\begin{enumerate}
\item $P$ has a root on $[a,r]$ if and only if $\Delta_{a,r}$ divides $N_P$.
\item In the remainder \eqref{eq:sphereremainder}, if $A=0$, divisibility forces $B=0$ and
      the entire sphere $[a,r]$ consists of roots.
\item If $A\ne0$, there is exactly one root on that sphere, namely $q=-BA^{-1}$.
\end{enumerate}
Real roots are exactly the real roots of $N_P$.
\end{theorem}
\begin{proof}
On $[a,r]$, evaluation is $P(q)=qA+B$, proving the alternatives once existence is known.
Modulo $\Delta_{a,r}$, the normal polynomial is congruent to
\[
 N(XA+B)=N(A)X^2+2\Rea(A\bar B)X+N(B).
\]
If $\Delta_{a,r}\mid N_P$ and $A=0$, this constant remainder forces $N(B)=0$, hence $B=0$.
If $A\ne0$, the degree-two polynomial must equal $N(A)\Delta_{a,r}$. Thus
\[
 \Rea(A\bar B)=-aN(A),\qquad N(B)=(a^2+r^2)N(A).
\]
For $q=-BA^{-1}$ these equations give $\Rea(q)=a$ and $N(q)=a^2+r^2$, so $q\in[a,r]$.
It is the required root. Conversely, if $qA+B=0$ at a point of the sphere and $A\ne0$,
substitute $B=-qA$ in the same two coefficients to obtain
$N(XA+B)=N(A)\Delta_{a,r}$. If $A=0$, then $B=0$. In both cases divisibility follows.
For real $x$, the final claim is $N_P(x)=N(P(x))$ and positivity of the norm.
\end{proof}

\begin{corollary}[Fundamental theorem and factorization]\label{cor:fta}
Every nonconstant one-sided polynomial over $\HF$ has a root. It admits a factorization
into $n$ linear factors and a final nonzero leading coefficient, with the displayed order of
factors significant. Its evaluation map $\HF\to\HF$ is surjective.
\end{corollary}
\begin{proof}
Over a real closed field, $N_P$ factors into real linear and irreducible quadratic factors.
A linear factor gives a real root. Any irreducible quadratic can be normalized to a
$\Delta_{a,r}$ and gives a root by \cref{thm:polynomialzeros}.
Use \eqref{eq:leftdivision} and induction for factorization. Apply the existence result to
$P-b$ for each target value $b$ to obtain surjectivity.
\end{proof}

This argument is a direct real-closed-field form of classical quaternionic zero theory;
compare the slice-algebra treatment in \cite{GP}. It is not an appeal to ordinary topological
degree on a compact sphere, which would need a different justification over $\No$.

\subsection{A precise degree count}
For a nonreal zero class $C=[a,r]$ define
$m_C=\ord_{\Delta_{a,r}}N_P$. For a real root $\lambda$ define
$m_\lambda=\tfrac12\ord_{X-\lambda}N_P$. The latter is an integer: repeatedly dividing
$P$ by the central factor $X-\lambda$ removes a factor $(X-\lambda)^2$ from its normal,
and the remaining normal is nonzero at $\lambda$.

\begin{proposition}[Conservation of class multiplicity]\label{prop:rootcount}
For a degree-$n$ polynomial,
\[
 \sum_{C\text{ nonreal zero class}}m_C+
 \sum_{\lambda\text{ real root}}m_\lambda=n.
\]
A spherical zero class has $m_C\ge2$. Consequently, if $s$ is the number of spherical
classes and $\ell$ the number of remaining individual zeros, including real zeros, then
\begin{equation}\label{eq:rootcount}
 2s+\ell\le n.
\end{equation}
\end{proposition}
\begin{proof}
Every irreducible factor of $N_P$ corresponds to a zero by \cref{thm:polynomialzeros}.
Its degree is $2n$, and the indicated multiplicities count half its factor degrees.
If the entire sphere vanishes, the remainder criterion gives $\Delta_{a,r}\mid P$.
Normal multiplicativity then gives $\Delta_{a,r}^2\mid N_P$. Every other zero class has
multiplicity at least one. These facts imply \eqref{eq:rootcount}.
\end{proof}

This is an effective algebraic algorithm relative to exact real-closed-field operations:
form $N_P$, factor it over $F$, compute the two-coefficient remainders, and distinguish
spherical from isolated roots. ``Effective relative to'' is essential: arbitrary surreal
coefficients do not come with finite computable encodings.

\begin{example}[Factor roots are not a list of evaluated roots]\label{ex:poly}
Consider
\[
 P=(X-\ii)\star(X-\jj)=X^2-X(\ii+\jj)+\kk.
\]
Then $P(\ii)=0$ but $P(\jj)=2\kk$. Its normal is $(X^2+1)^2$.
Modulo $X^2+1$ the remainder is $-X(\ii+\jj)+(\kk-1)$, so the unique root on the only
possible class is $\ii$. This isolated root has class multiplicity two.
The two written linear factors therefore do not supply two distinct pointwise zeros.
\end{example}

Central polynomials behave differently: a nonreal root forces its entire conjugacy sphere
to vanish. Conversely no theorem here says that every arbitrary noncommutative polynomial
expression has a root. For example $\ii q-q\ii=1$ has no solution because its left side has
zero scalar part. The one-sided, central-indeterminate convention is a mathematical
hypothesis, not merely notation.

\section{Quaternionic linear algebra and multiscale spectral stability}\label{sec:spectral}

On the right vector space $\HF^n$ use
\[
 \angles{x,y}=\sum_{j=1}^n\bar x_jy_j,
 \qquad |x|=\sqrt{\angles{x,x}}.
\]
Matrices act on the left. The adjoint is $A^*=\bar A^{\,T}$, and a matrix is unitary when
$U^*U=I$. Orthogonalization uses finite sums and positive square roots, so works over a real
closed field exactly as a finite algebraic procedure.

\begin{theorem}[Hermitian spectral theorem]\label{thm:spectral}
If $A=A^*\in M_n(\HF)$, there is a unitary $U$ and real scalars
$\lambda_1,\ldots,\lambda_n\in F$ such that
\[
 U^*AU=\diag(\lambda_1,\ldots,\lambda_n).
\]
Every positive definite Hermitian matrix has a unique positive definite Hermitian square
root. Every quaternion matrix admits a singular value decomposition; in particular an
invertible square matrix has a polar decomposition.
\end{theorem}
\begin{proof}
Write $A=Z+W\jj$ with complex matrices over $C_F$. Its complex representation is
\[
 \rho(A)=\begin{pmatrix}Z&W\\-\bar W&\bar Z\end{pmatrix}.
\]
The identification $x+y\jj\mapsto(x,-\bar y)^T$ intertwines its action with that of $A$.
If $A$ is Hermitian, so is $\rho(A)$. Since $C_F$ is algebraically closed, the latter has an
eigenvector $v\ne0$ with eigenvalue $\lambda\in C_F$. The scalar $v^*v$ is positive in $F$,
and $v^*\rho(A)v$ is fixed by conjugation. Thus
$\lambda=(v^*\rho(A)v)/(v^*v)\in F$.
The corresponding quaternion vector is a right eigenvector with the same real eigenvalue.
Normalize it. Its orthogonal complement is invariant under $A$, because $A=A^*$ and the
eigenvalue is central. Induction on $n$ gives the diagonalization.

For a positive definite matrix replace each positive diagonal eigenvalue by its positive
square root. Uniqueness follows because a positive square root commutes with its square,
preserves its eigenspaces, and is positive scalar on each after diagonalization there.
Diagonalize $A^*A$ for singular values; for each nonzero singular value $\sigma$, normalize
$Av$ by $\sigma$ and complete orthonormal bases for the remaining spaces. This proves the
singular value decomposition. For invertible $A$, writing
$A=(A(A^*A)^{-1/2})(A^*A)^{1/2}$ gives the polar factors.
\end{proof}

No appeal to compactness of a surreal unit sphere occurs in this proof. It uses only finite
linear algebra and algebraic closedness of the complexification of the center.

Define $\norm A$ to be its largest singular value. The singular value decomposition proves
$|Ax|\le\norm A|x|$, with equality for a suitable unit vector. For a square $n\times n$
matrix,
\begin{equation}\label{eq:matrixnorm}
 \max_{i,j}|a_{ij}|\le\norm A\le n\max_{i,j}|a_{ij}|.
\end{equation}
Consequently for a nonzero surquaternion matrix,
\begin{equation}\label{eq:matrixvaluation}
 w(\norm A)=\min_{i,j}w(a_{ij}).
\end{equation}
Matrix multiplication may cancel leading terms, so only
$w(AB)\ge w(A)+w(B)$ is automatic; the equality for nonzero individual quaternions must
not be copied to matrices.

\subsection{Weyl's estimate without compactness}
Order the eigenvalues of a Hermitian matrix increasingly. Finite-dimensional spectral
algebra gives the min--max formula
\[
 \lambda_k(A)=\min_{\dim W=k}\ \max_{x\in W,\ |x|=1}\angles{x,Ax}.
\]
Indeed the span of the first $k$ eigenvectors realizes the upper bound. Any $k$-dimensional
right subspace meets the span of eigenvectors $k,\ldots,n$ nontrivially, proving the lower
bound. The inner maximum exists by applying the spectral theorem to the compressed
Hermitian form on $W$. Thus this notation represents attained extrema justified by linear
algebra, not by a general extreme-value theorem over $F$.

Since $|\angles{x,Ex}|\le\norm E$ for unit $x$, it follows that
\begin{equation}\label{eq:weyl}
 |\lambda_k(A+E)-\lambda_k(A)|\le\norm E
 \qquad(A=A^*,\ E=E^*).
\end{equation}
If $A$ has finite entries, its eigenvalues are finite by \eqref{eq:matrixnorm}; reducing a
unitary diagonalization by standard part proves that their ordered standard parts are exactly
the ordered eigenvalues of $\st(A)$, including multiplicities.

\begin{theorem}[A spectral-projector estimate at arbitrary surreal gap scale]\label{thm:projector}
Let $A=A^*\in M_n(\HF)$, $n\ge2$, have a simple eigenvalue $\lambda$. Let
\[
 g=\min_{\mu\ne\lambda}|\mu-\lambda|>0,
 \qquad e=\norm E<g/2,
\]
where $E=E^*$. There is exactly one eigenvalue $\lambda'$ of $A+E$ in
$(\lambda-g/2,\lambda+g/2)$, counted with quaternionic multiplicity.
If $P,P'$ are the associated rank-one orthogonal projections, then
\begin{equation}\label{eq:projectorbound}
 \norm{P'-P}\le\frac{e}{g-e}.
\end{equation}
For surreal coefficients, the sufficient condition $w(E)>w(g)$ gives
\begin{equation}\label{eq:projectorval}
 w(P'-P)\ge w(E)-w(g),
\end{equation}
with zero matrices interpreted as having valuation $+\infty$.
\end{theorem}
\begin{proof}
Weyl's estimate places the matched ordered eigenvalue within $e$ of $\lambda$, and every
other perturbed eigenvalue at least $g-e>g/2$ from $\lambda$. This proves uniqueness and
simplicity. Let $v'$ be a unit vector for $\lambda'$, and write
$v'=Pv'+y$, $y=(I-P)v'$. Applying $I-P$ to
$(A+E-\lambda')v'=0$ gives
\[
 (A-\lambda')y=-(I-P)Ev'.
\]
On $\operatorname{im}(I-P)$, every spectral value of $A-\lambda'$ has absolute value at
least $g-|\lambda'-\lambda|\ge g-e$. The spectral theorem therefore gives
$|y|\le e/(g-e)$.

For two rank-one orthogonal projections, $\norm{P'-P}=|(I-P)v'|$.
To verify this over $\HF$, multiply one generating vector by a unit quaternion so that its
inner product with the other is a nonnegative scalar, and restrict to their two-dimensional
span. The resulting two-by-two real symmetric matrix has eigenvalues
$\pm\sqrt{1-|\angles{v,v'}|^2}$; the orthogonal complement contributes zero.
This proves \eqref{eq:projectorbound}.
If $w(E)>w(g)$, then $e/g$ is infinitesimal, so $e<g/2$ and $w(g-e)=w(g)$.
Taking valuations and using \eqref{eq:matrixvaluation} yields \eqref{eq:projectorval}.
\end{proof}

The estimate is useful when the standard-part matrix has a multiple eigenvalue: an
infinitesimal gap can still resolve its eigenvectors at a finer scale. Replacing $g$ by an
ordinary standard-part gap would discard exactly that information.

\begin{example}[A sharp loss of gap scale]\label{ex:spectral}
Let $t=\omega^{-1}$, $0<\alpha<\beta$, and
\[
 A=\begin{pmatrix}0&t^\beta\ii\\-t^\beta\ii&t^\alpha\end{pmatrix}.
\]
This is a perturbation of $\diag(0,t^\alpha)$ of norm $t^\beta$, with gap $t^\alpha$.
Its eigenvalues are exactly
\begin{equation}\label{eq:twoeigen}
 \lambda_\pm=\frac{t^\alpha\pm\sqrt{t^{2\alpha}+4t^{2\beta}}}{2}.
\end{equation}
The positive-order binomial expansion gives
\begin{align*}
 \lambda_-&=-t^{2\beta-\alpha}+t^{4\beta-3\alpha}
             +\text{terms of valuation at least }6\beta-5\alpha,\\
 \lambda_+&=t^\alpha+t^{2\beta-\alpha}-t^{4\beta-3\alpha}
             +\text{terms of valuation at least }6\beta-5\alpha.
\end{align*}
The lower spectral projector is
$P_-=(\lambda_+I-A)/(\lambda_+-\lambda_-)$.
Its off-diagonal entry is
\[
 -\frac{t^\beta\ii}{\sqrt{t^{2\alpha}+4t^{2\beta}}}
 =-t^{\beta-\alpha}\ii+\text{terms of strictly higher valuation}.
\]
Its diagonal changes have valuation $2(\beta-\alpha)$. Thus
$w(P_--\diag(1,0))=\beta-\alpha$, realizing the loss in \eqref{eq:projectorval} exactly.
The eigenvalue displacement occurs at the finer scale $2\beta-\alpha$.
\end{example}

\section{Infinitesimal Lie theory and quaternionic exponentials}\label{sec:explog}

\subsection{The canonical neighborhood of the identity}
For $X,Y\in\mm_H$ define
\begin{equation}\label{eq:explog}
 \Exp_{\mm}(X)=\sum_{n\ge0}\frac{X^n}{n!},\qquad
 \Log_{\mm}(1+Y)=\sum_{n\ge1}\frac{(-1)^{n+1}Y^n}{n}.
\end{equation}
Both families are strongly summable by the positive-support lemma. The constant coefficient
is handled as a single term, not as an ordinary convergent sum of constants.

\begin{theorem}[Principal logarithm and infinitesimal rotations]\label{thm:principal-log}
The maps in \eqref{eq:explog} are mutually inverse bijections between $\mm_H$ and
$1+\mm_H$. They commute with conjugation and with inner conjugation. Their restrictions give
a bijection
\begin{equation}\label{eq:logunits}
 \Exp_{\mm}:\mm_H\cap\Ima\SQ\ \longleftrightarrow\ K:
 \Log_{\mm},
\end{equation}
where $K$ is the kernel in \eqref{eq:unitsreduction}. For commuting infinitesimals,
$\Exp_{\mm}(X+Y)=\Exp_{\mm}(X)\Exp_{\mm}(Y)$.
\end{theorem}
\begin{proof}
The one-variable formal identities $\log(\exp T)=T$ and
$\exp(\log(1+T))=1+T$ can be substituted at positive valuation. At each exponent only
finitely many words contribute, so formal regrouping becomes valid Hahn regrouping.
Each identity involves powers of one element, which generate a commutative subalgebra;
noncommutativity elsewhere is irrelevant. Coefficientwise conjugation and inner conjugation
preserve these identities and strong supports.

If $X$ is imaginary, $\overline{\Exp_{\mm}(X)}=\Exp_{\mm}(-X)$, hence the norm is $1$.
Conversely, if $u\in K$, then $\bar u=u^{-1}$. Therefore
$\overline{\Log_{\mm}u}=\Log_{\mm}(u^{-1})=-\Log_{\mm}u$, so its logarithm is imaginary.
The commuting addition law follows from the formal binomial expansion and the same support
regrouping.
\end{proof}

This is a bijection, not an assertion that the imaginary additive group is the group of
infinitesimal rotations. The transported group law is noncommutative.

\begin{proposition}[Baker--Campbell--Hausdorff at positive valuation]\label{prop:bch}
For $X,Y\in\mm_H$ the expression
\[
 \BCH(X,Y)=\Log_{\mm}\bigl(\Exp_{\mm}(X)\Exp_{\mm}(Y)\bigr)
\]
is defined by strong Hahn summation. Its initial terms are
\begin{align}\label{eq:bch}
 \BCH(X,Y)={}&X+Y+\tfrac12[X,Y]
 +\tfrac1{12}[X,[X,Y]]+\tfrac1{12}[Y,[Y,X]]+R_4.
\end{align}
If $w(X),w(Y)\ge\delta>0$, then $w(R_4)\ge4\delta$.
Writing $\alpha=w(X)$, $\beta=w(Y)$ for nonzero $X,Y$, the group commutator satisfies
\begin{equation}\label{eq:groupcomm}
 \Exp X\Exp Y\Exp(-X)\Exp(-Y)
 =1+[X,Y]+R,
 \quad w(R)\ge\alpha+\beta+\min\{\alpha,\beta\}.
\end{equation}
Here each exponential is the infinitesimal one.
\end{proposition}
\begin{proof}
Use the formal exponential and logarithm in the free associative power-series algebra on
two variables, filtered by total word degree. Their compositions have finitely many words
at each degree, and expansion through degree three gives \eqref{eq:bch}. At $X,Y$ with
positive valuation, every word is controlled by the monoid generated by their combined
positive supports. The support lemma makes evaluation and regrouping legitimate across
all degrees. Terms of degree at least four have valuation at least $4\delta$.
In the group commutator, the degree-two term is $XY-YX$. Every remaining nonzero word
involves both variables and at least one additional occurrence, giving the stated bound.
\end{proof}

For imaginary $X,Y$, $[X,Y]=2X\times Y$. If their leading imaginary quaternion coefficients
are not parallel, the commutator has valuation $\alpha+\beta$ and is nonzero. For example
$X=t\ii$, $Y=t^2\jj$ give
\[
 \Exp X\Exp Y\Exp(-X)\Exp(-Y)=1+2t^3\kk+
 \text{terms of valuation at least }4.
\]
Thus two infinitesimal rotations can agree with the identity to distinct orders but produce a
third, smaller rotational scale when commuted. The associated graded Lie bracket is precisely
twice the ordinary cross product, accompanied by addition of exponents.
The bounds do not assume $n\delta$ cofinal and do not turn the Hahn identities into limits
of a sequence in the full fine topology.

\subsection{Canonical finite angles}
For a finite surreal $\theta=r+\eps$, $r\in\R$ and $\eps$ infinitesimal, define
\begin{align}\label{eq:finite-trig}
 \sin\theta&=\sin r\sum_{n\ge0}\frac{(-1)^n\eps^{2n}}{(2n)!}
            +\cos r\sum_{n\ge0}\frac{(-1)^n\eps^{2n+1}}{(2n+1)!},\\
 \cos\theta&=\cos r\sum_{n\ge0}\frac{(-1)^n\eps^{2n}}{(2n)!}
            -\sin r\sum_{n\ge0}\frac{(-1)^n\eps^{2n+1}}{(2n+1)!}.
\end{align}
The ordinary real values are taken first; only the infinitesimal correction is Hahn-summed.
This is a restricted analytic extension, consistent with the established surreal analytic
structure \cite{vdDE,EK}. Formal addition formulas and the ordinary real addition formulas
prove their finite-surreal analogues.

\begin{proposition}[Every direction has a finite angle]\label{prop:finiteangles}
For $c,s\in\No$ with $c^2+s^2=1$ there exists a finite surreal $\theta$, unique modulo
$2\pi\Z$, such that $c+\ii s=\cos\theta+\ii\sin\theta$.
If $s\ge0$, there is a unique representative $\theta\in[0,\pi]$.
\end{proposition}
\begin{proof}
The coordinates are finite and their standard parts lie on the ordinary unit circle. Choose
an ordinary angle $r$ of that standard part. Multiplication by $\cos r-\ii\sin r$ leaves a
unit in $1+\mm_H$, within the commutative slice $\No_{\ii}$. Its principal logarithm is
$\ii\eps$ for a unique infinitesimal $\eps$, by \cref{thm:principal-log}. Then
$\theta=r+\eps$ gives the desired point. If two angles give it, their standard-part
difference is $2\pi k$ for an ordinary integer $k$; injectivity of the principal logarithm
forces the remaining infinitesimal difference to vanish. On the upper semicircle choose the
ordinary representative in $[0,\pi]$. At its endpoints, the sign of the infinitesimal sine
correction selects the inward direction, so the corrected angle stays in this interval.
\end{proof}

This proposition supplies angles even for directions determined by vectors of infinite or
infinitesimal length. For $u=a+v\in S^3(\No)$ write $a=\cos\theta$ and
$|v|=\sin\theta$ with $\theta\in[0,\pi]$. If $v\ne0$, set $I=v/|v|$; then
$u=\cos\theta+I\sin\theta$. For $u=\pm1$, any $I$ is available. Formula
\eqref{eq:rodrigues} becomes rotation about $I$ through the finite angle $2\theta$.

\subsection{An exponential with unrestricted real part}
Let
\[
 \mathcal T=\{a+v\in\SQ:a\in\No,\ v\in\Ima\SQ,\ |v|\text{ finite}\}.
\]
Using Gonshor's real exponential and the finite-angle trigonometry, define
\begin{equation}\label{eq:tubeexp}
 \Exp_0(a+v)=
 \begin{cases}
 e^a\left(\cos r+\dfrac{v}{r}\sin r\right),&r=|v|>0,\\
 e^a,&v=0.
 \end{cases}
\end{equation}
This definition is canonical within the chosen standard surreal exponential and restricted
analytic structure. Its domain contains every finite quaternion and all central surreals,
but not $\omega\ii$.

\begin{theorem}[Polar logarithms and arbitrary-scale roots]\label{thm:polarlog}
The map $\Exp_0$ is conjugation-equivariant, satisfies
$|\Exp_0(a+v)|=e^a$, and obeys the addition law on commuting pairs in $\mathcal T$.
It is surjective onto $\SQ^\times$. Specifically, if
\[
 q=\rho(\cos\theta+I\sin\theta),\qquad
 \rho=|q|>0,\quad0\le\theta\le\pi,
\]
then
\begin{equation}\label{eq:polarlog}
 \Exp_0(\log\rho+I\theta)=q.
\end{equation}
For every ordinary integer $n\ge1$, one $n$th root is
\begin{equation}\label{eq:nthroot}
 \rho^{1/n}\left(\cos(\theta/n)+I\sin(\theta/n)\right).
\end{equation}
\end{theorem}
\begin{proof}
Inner conjugation fixes the scalar part and the length of the imaginary part, and carries
its direction to the conjugate direction. The norm formula is $\cos^2r+\sin^2r=1$.
A commuting pair lies in a common slice. There \eqref{eq:tubeexp} can be written
$e^a(\cos b+I\sin b)$ with a signed finite coordinate $b$, and the addition law is the
finite sine and cosine addition formula. Every unit quaternion has the polar representation
just proved, yielding \eqref{eq:polarlog}. The addition formula, applied $n$ times in the
same slice, proves \eqref{eq:nthroot}.
\end{proof}

There is no single-valued logarithm choice implicit on the negative scalar axis: there every
$\log\rho+\pi I$ with $I\in\mathbb S(\No)$ exponentiates to $-\rho$. Even off that axis,
ordinary $2\pi$ periods give other logarithms. Surjectivity of a quaternion exponential need
not entail a preferred logarithm.

On finite quaternions, \eqref{eq:tubeexp} agrees with the four-real-variable Taylor extension
of the ordinary quaternion exponential. At a nonzero ordinary imaginary radius this follows
by analytic composition with the square root; at radius zero the removable function
$\sin r/r$ is evaluated by its even power series. If $q=q_0+\eps$ with an ordinary quaternion
$q_0$ and infinitesimal $\eps$, this is \emph{not} generally
$\exp(q_0)\Exp_{\mm}(\eps)$: the two summands need not commute.

Nor is \eqref{eq:tubeexp} obtained by strongly summing $\sum q^n/n!$ for every $q\in\mathcal T$.
At a nonzero ordinary scalar, all those terms have exponent zero, violating finite overlap.
At $q=\omega\ii$ their exponents are unbounded in the wrong direction. These are distinct
failures, neither repaired by citing ordinary real convergence.

\subsection{A global exponential from Ehrlich--Kaplan trigonometry}\label{sec:globalexp}
There is established global surreal trigonometry; it must not be dismissed as impossible.
Ehrlich and Kaplan specify it using the initial integer part $\Oz$ \cite[Section 11]{EK}.
In Conway normal forms, $\Oz$ consists of a purely infinite part and an ordinary integer
constant term. For each $x\in\No$, take the unique $d\in\Oz$ for which
\[
 r=x-2\pi d\in[0,2\pi),
\]
and set $\sin_{\mathrm{EK}}x=\sin r$, $\cos_{\mathrm{EK}}x=\cos r$, with the restricted
values on this bounded interval. The resulting circle character has period group
$2\pi\Oz$, agrees with \eqref{eq:finite-trig} on finite inputs, and gives a surcomplex
exponential. This is a specified global normalization, not something forced by the bare
ordered-field structure alone.

It also gives a global surquaternion exponential:
\begin{equation}\label{eq:globalexp}
 \Exp_{\mathrm{EK}}(a+v)
 =e^a\left(\cos_{\mathrm{EK}}r+\frac{v}{r}\sin_{\mathrm{EK}}r\right)
 \quad(r=|v|>0),
\end{equation}
with $e^a$ when $v=0$. The proof of \cref{thm:polarlog} still shows conjugation equivariance,
the norm formula, surjectivity, and the addition law for commuting pairs. A noncommuting
addition law is neither expected nor asserted. On each fixed complex slice the usual
exponential homomorphism has kernel $2\pi\Oz I$; on the entire noncommutative algebra
$\Exp_{\mathrm{EK}}$ is not a homomorphism of additive groups, so calling its full preimage
of $1$ a group kernel would be misleading.

The global choice has striking large-scale consequences. Since $\omega/(2\pi)\in\Oz$,
\begin{equation}\label{eq:omegaphase}
 \sin_{\mathrm{EK}}\omega=0,\qquad
 \cos_{\mathrm{EK}}\omega=1,\qquad
 \Exp_{\mathrm{EK}}(\omega\ii)=1.
\end{equation}
This does not identify a divergent or nonsummable Taylor series with $1$. It computes a value
of a separately defined global operation. The next section shows why differentiating this
operation with a standard surreal derivation requires special care.

\section{Derivations and an obstruction to oscillatory calculus}\label{sec:derivations}

\subsection{All algebraic derivations have a scalar and an inner part}
Let $d:F\to F$ be a field derivation. There is a componentwise extension
\begin{equation}\label{eq:coeffderiv}
 D_0(a+b\ii+c\jj+d_1\kk)
 =d(a)+d(b)\ii+d(c)\jj+d(d_1)\kk.
\end{equation}
It satisfies the Leibniz rule and commutes with conjugation because Hamilton's structure
constants are integers.

\begin{theorem}[Derivation decomposition]\label{thm:derivations}
Every ring derivation $D:\HF\to\HF$ can be written uniquely as
\begin{equation}\label{eq:derivdecomp}
 D=D_0+\ad_b,\qquad \ad_b(q)=bq-qb,
\end{equation}
where $D_0$ extends its restriction $d$ to the center componentwise and $b\in\Ima\HF$.
\end{theorem}
\begin{proof}
Differentiating the commutation relation $cq=qc$ for $c\in F$ shows that $D(c)$ commutes
with every quaternion. Thus $D$ restricts to a derivation of $F$. Subtracting its
componentwise extension gives an $F$-linear derivation $\delta$.

For an imaginary quaternion $u$, differentiate $u^2=-N(u)$. Since $\delta$ vanishes on $F$,
this implies that $\delta(u)$ has zero scalar part and is orthogonal to $u$ (first do this
for nonzero $u$). Applying the result to $u+v$ shows that $\delta|_V$ is skew-symmetric.
Every skew-symmetric three-by-three matrix is uniquely $u\mapsto2b\times u$ for a vector
$b\in V$, by its three independent entries. This is exactly $[b,u]$; both maps vanish on
$F$, proving \eqref{eq:derivdecomp}. Equivalently, the imaginary representative is explicitly
\[
 b=-\frac14\bigl(\delta(\ii)\ii+\delta(\jj)\jj+\delta(\kk)\kk\bigr).
\]
Uniqueness also follows because an imaginary element commuting with the whole algebra is
zero.
\end{proof}

For the componentwise extension one has
\begin{align}\label{eq:derivformulas}
 D_0(q^{-1})&=-q^{-1}D_0(q)q^{-1},\\
 D_0N(q)&=2\Rea(\bar qD_0q),\\
 D_0(q^n)&=\sum_{k=0}^{n-1}q^k(D_0q)q^{n-1-k}.
\end{align}
The last formula becomes $nq^{n-1}D_0q$ only when the relevant factors commute. The inner
part in \eqref{eq:derivdecomp} gives an algebraic analogue of a change of quaternionic frame;
its commutator cannot be absorbed into a scalar derivative of the coefficients.

\subsection{The Berarducci--Mantova extension}
Fix the Berarducci--Mantova derivation $\partial$ on $\No$ normalized by
$\partial\omega=1$. We use the established properties that it is strongly additive,
commutes with the real surreal exponential, has constant field $\R$, is surjective, and
maps infinitesimals to infinitesimals \cite[Theorems A and B]{BM}. Let $D$ denote its componentwise extension
to $\SQ$, with $D\ii=D\jj=D\kk=0$.

It follows immediately that $\ker D=\HR$, that $D$ is surjective, and that it respects
strongly summable families componentwise. For an infinitesimal $X$, differentiating its
strong exponential gives the valid noncommutative formula
\begin{equation}\label{eq:dexp}
 D\Exp_{\mm}(X)=
 \sum_{n\ge1}\frac1{n!}\sum_{k=0}^{n-1}X^k(DX)X^{n-1-k}.
\end{equation}
Positive supports make this expression strongly summable. If $[X,DX]=0$, it reduces to
$\Exp_{\mm}(X)DX$. Nothing in this local assertion implies a global commuting chain rule
for every possible extension of the quaternion exponential.

\begin{lemma}[Finite inputs have infinitesimal derivatives]\label{lem:small}
If $q\in\OO_H$, then $Dq\in\mm_H$.
\end{lemma}
\begin{proof}
Write $q=\st(q)+\eps$. The standard part belongs to $\HR$ and is killed by $D$.
Each coordinate of $\eps$ is infinitesimal and its derivative is infinitesimal by the smallness
property of $\partial$. The four-coordinate assertion follows.
\end{proof}

\begin{theorem}[No constant-frequency quaternionic oscillation]\label{thm:nooscillation}
Let $I\in\mathbb S(\R)$ and $c\in\R\setminus\{0\}$. Neither equation
\begin{equation}\label{eq:oscillation}
 Dy=cIy\qquad\text{or}\qquad Dy=ycI
\end{equation}
has a nonzero solution in $\SQ$. In particular, $Dy=\ii y$ has only the zero solution, even
in a fixed surcomplex slice.
\end{theorem}
\begin{proof}
For the left equation,
\[
 D N(y)=\overline{Dy}\,y+\bar y\,Dy
       =-\bar y\,cI y+\bar y\,cI y=0.
\]
The right equation gives the same cancellation, using centrality of $N(y)$.
Thus $N(y)$ lies in the constant field $\R$. If $y\ne0$, this is a positive ordinary real
number, and $|y|=\sqrt{N(y)}$ is a positive ordinary real as well. Every coordinate of $y$
is therefore finite. By \cref{lem:small}, $Dy$ has infinitesimal coordinates, so its norm is
infinitesimal or zero. But the right side of either equation has norm $|c|\,|y|$, a positive
ordinary real. This is impossible.
\end{proof}

The proof uses smallness and the constant field, not any assumption about the global
Ehrlich--Kaplan phase. It therefore rules out an entire compatibility requirement, rather
than merely exposing a defect of one phase choice.

\begin{corollary}[No global commuting chain rule]\label{cor:nochain}
There is no map $E:\SQ\to\SQ^\times$ satisfying
\[
 D(E(q))=E(q)Dq
\]
for every $q$ for which $q$ and $Dq$ commute. The same conclusion holds with the factors
on the right side reversed. In particular the usual commuting exponential chain rule cannot
hold at $q=\omega\ii$.
\end{corollary}
\begin{proof}
At $q=\omega\ii$ we have $Dq=\ii$ and $[q,Dq]=0$. A nonzero $E(q)$ would solve one of the
forbidden equations in \cref{thm:nooscillation}.
\end{proof}

For the actual global normalization in \eqref{eq:globalexp}, the contradiction is already
visible numerically as an exact symbolic identity:
\[
 D\Exp_{\mathrm{EK}}(\omega\ii)=D1=0,
 \qquad \Exp_{\mathrm{EK}}(\omega\ii)D(\omega\ii)=\ii.
\]
The normalization is still a valid global algebraic exponential on each commuting slice.
It simply does not have this differential compatibility.

These conclusions do not conflict with the existence of real antiderivatives or the
Liouville-closed real differential-field properties proved in \cite{BM}. A constant
imaginary frequency is not a scalar coefficient in that real differential field. Adjoining
formal oscillations would require a larger differential algebra or changing the derivation;
ordinary quaternion scalar extension by itself does not produce them.

Theorems \ref{thm:nooscillation} and \ref{cor:nochain} are proved here from the cited
properties of the derivation. We make no claim that this formulation has never appeared in
the literature. Their role is to identify a concrete boundary that any proposed unified
surquaternion calculus must respect.

\section{A support-controlled analytic theory}\label{sec:analysis}

The algebra has many complex slices, but a satisfactory analytic theory must relate them
without confusing three operators: differentiation of a polynomial or analytic \emph{variable},
the componentwise field derivation $D$ in \cref{sec:derivations}, and the four-dimensional
Fueter operator. They have different domains and different defining properties.

This section constructs a controlled extension of classical slice analysis \cite{GP}.
Every infinite expression uses strong supports or ordinary coefficient analysis. No assertion
about arbitrary class functions on $\SQ$ is intended.

\subsection{Formal functions on an infinitesimal disk}
A simple useful class consists of right-coefficient series
\begin{equation}\label{eq:formalseries}
 f(x)=\sum_{n\ge0}x^na_n,
 \qquad a_n\in D_\Gamma,
\end{equation}
for which all $\supp(a_n)$ lie in one well-ordered set $S\subset\Gamma$.

\begin{proposition}[Uniform-support evaluation]\label{prop:evaluation}
Under this common-support hypothesis, \eqref{eq:formalseries} is strongly summable at every
$x\in\mm_H$ in a common workspace. Its formal derivative
$\sum_{n\ge1}nx^{n-1}a_n$ has the same property. Products and formal differentiation are
compatible with strong regrouping whenever the same common-support conditions are imposed
on the participating series.
\end{proposition}
\begin{proof}
Let $T$ be the positive support of $x$. Every summand is supported in $S+nT$. The monoid
$T^*$ is well ordered by the support lemma, and $S+T^*$ is well ordered. For a fixed exponent,
there are only finitely many pairs consisting of an exponent of $S$ and one of $T^*$, and
only finitely many words representing the latter, across all lengths. Thus only finitely many
summands contribute. Shifting the index by one proves the derivative assertion. The same
word-count argument proves the asserted regroupings and the formal product rule.
\end{proof}

The derivative here is the \emph{slice/formal} derivative, not a claim that
$df_x(h)=hf'(x)$ for arbitrary noncommuting $h$. Even the elementary polynomial $x^2$ has
full algebraic differential $h\mapsto xh+hx$, rather than $h\mapsto2hx$ in all directions.

The common-support hypothesis is substantive. In the workspace with $\Gamma=\Q$, the
series
\[
 \sum_{n\ge0}t^{-n^2}x^n
\]
is not strongly summable at any nonzero argument: if $w(x)=\delta\in\Q$, the leading
exponents $-n^2+n\delta$ are unbounded below. Enlarging the workspace can change this:
at the surreal argument $x=t^\omega$ these leading exponents increase and the monomial
family is strongly summable. A failure in a fixed value group is not automatically a failure
at every surreal scale.

A central translation and scaling replace $x$ by $(q-a)/r$ with $a,r\in\No$, $r\ne0$.
This gives controlled infinitesimal disks even around an infinite scalar center. It does not
supply arbitrary expansions about noncentral centers: the order of increments and center
terms then matters and must be retained in noncommutative words.

\subsection{Common-domain Hahn stems}
Let $\Omega\subset\C$ be an ordinary connected conjugation-invariant open set meeting the
real axis. Introduce an \emph{auxiliary central} complex unit $\iota$, distinct from every
quaternion unit. An ordinary stem coefficient is a holomorphic map
\[
 F_\gamma:\Omega\to\HR\otimes_\R\C
 \quad\text{with}\quad
 F_\gamma(\bar z)=\overline{F_\gamma(z)}^{\,\iota},
\]
where the overline changes $\iota$ to $-\iota$ and leaves ordinary quaternion coefficients
unchanged. Holomorphic means holomorphic in its four complex scalar coordinates.
The target is a complexified quaternion algebra, not the division algebra $\SQ$.

A \emph{common-domain Hahn stem} is
\begin{equation}\label{eq:hahnstem}
 F(z)=\sum_{\gamma\in S}t^\gamma F_\gamma(z),
\end{equation}
with $S$ a well-ordered set and the same $\Omega$ for every coefficient. At a finite
$z=z_0+\eta\in\No[\iota]$, $z_0=\st(z)\in\Omega$, evaluate by
\begin{equation}\label{eq:stemeval}
 F(z)=\sum_{\gamma\in S}\sum_{n\ge0}
 t^\gamma\frac{F_\gamma^{(n)}(z_0)}{n!}\eta^n.
\end{equation}
The factors $\eta$ are central in this complexification. The support proof of
\cref{prop:evaluation} makes the double family strongly summable. It requires no common
bound on the magnitudes of all ordinary derivatives; finiteness at each exponent replaces
an analytic sum over $\gamma$.

Write the result $F(a+\iota b)=A(a,b)+\iota B(a,b)$ with $A,B\in\SQ$. Its induced
\emph{slice function} is
\begin{equation}\label{eq:slicefunction}
 f(a+Ib)=A(a,b)+I B(a,b),\qquad I\in\mathbb S(\No),
\end{equation}
for finite $a,b$ with $\st(a+\iota b)\in\Omega$. The parity properties
$A(a,-b)=A(a,b)$, $B(a,-b)=-B(a,b)$ make this independent of the representation using
$(-I,-b)$; at $b=0$, $B=0$ makes it independent of $I$ altogether. Values may be infinite,
since $S$ can contain negative exponents, even though the argument domain in this
construction is a finite standard-part thickening.

The Cauchy--Riemann identities for the stem are
\begin{equation}\label{eq:CR}
 A_a=B_b,\qquad A_b=-B_a.
\end{equation}
Consequently on a fixed slice,
$\tfrac12(\partial_a+I\partial_b)f=0$. These variable derivatives mean ordinary coefficient
derivatives followed by the same Taylor evaluation. They are not the operator $D$ acting
on the surreal coefficients as numbers.

\begin{theorem}[Representation and a coherent identity principle]\label{thm:representation}
For every slice function of the form \eqref{eq:slicefunction},
\begin{equation}\label{eq:representation}
 f(a+Jb)=\frac12\bigl[(1-JI)f(a+Ib)+(1+JI)f(a-Ib)\bigr]
\end{equation}
for $I,J\in\mathbb S(\No)$. If a common-domain Hahn slice function vanishes on an ordinary
nonempty real interval in $\Omega$, then it vanishes throughout its induced domain.
\end{theorem}
\begin{proof}
Substitute $A+IB$ and $A-IB$ on the right of \eqref{eq:representation}; the terms simplify
to $A+JB$. No coefficient has been moved across $I$ or $J$.
At an ordinary real point $s$, the function is
$\sum_\gamma t^\gamma F_\gamma(s)$ with ordinary quaternion coefficients. If it is zero,
uniqueness of Hahn coefficients makes every $F_\gamma(s)$ zero. Vanishing on a real interval
and the ordinary complex identity theorem, applied to its four holomorphic coordinates,
make each $F_\gamma$ zero on connected $\Omega$. Evaluation then gives zero everywhere.
\end{proof}

This identity principle is deliberately coherent and coefficientwise. It does not claim that
an arbitrary fine-differentiable class function with many zeros is identically zero;
\cref{ex:zeroderivative} already demonstrates the danger of such a claim.

\subsection{Products, reciprocals, and zero spheres}
If two stems have values $A+\iota B$ and $C+\iota D_1$, their product induces
\begin{equation}\label{eq:starfunction}
 (f\star g)(a+Ib)=(AC-BD_1)+I(AD_1+BC).
\end{equation}
This is the slice product, and need not equal $(A+IB)(C+ID_1)$. Supports of the product
lie in the sum of two well-ordered sets, with finite coefficient decompositions, so common-domain
Hahn stems are closed under this product.
At a point where $f(q)\ne0$, the same twist as in \eqref{eq:twisted} holds:
\begin{equation}\label{eq:twistedfunction}
 (f\star g)(q)=f(q)g\bigl(f(q)^{-1}qf(q)\bigr).
\end{equation}
Indeed conjugating $q=a+Ib$ replaces $I$ by $I'=f(q)^{-1}If(q)$, while leaving $a,b$ fixed.
Then $f(q)(C+I'D_1)=f(q)C+If(q)D_1$, which expands to \eqref{eq:starfunction}.
If $f(q)=0$, substitution in that formula gives $(f\star g)(q)=0$ as well.

Quaternion-conjugating the coefficients of a stem defines $f^c$. Its normal
$f^s=f\star f^c=f^c\star f$ has scalar complex stem values:
\[
 (A+\iota B)(\bar A+\iota\bar B)
 =N(A)-N(B)+2\iota\Rea(A\bar B).
\]
Thus $f^s$ is \emph{slice-preserving}: at $a+Ib$ its value lies in $\No_I$, not necessarily
in the center $\No$. Wherever its normal is nonzero and a reciprocal stem is admitted
locally, the slice inverse evaluates as
\begin{equation}\label{eq:starinverse}
 f^{-*}(q)=f^s(q)^{-1}f^c(q).
\end{equation}
The order displayed is part of the formula. Local admission requires support control:
for example, on an ordinary neighborhood where a leading scalar coefficient of the normal
has no zeros, factor out that coefficient and monomial and invert the positive-order
remainder by \cref{prop:neumann}. A pointwise nonzero normal is not by itself a proof that
one global common-domain reciprocal family exists without shrinking or changing domains.

On one conjugacy sphere the value is $A+IB$. If $B=0$, either the entire sphere vanishes
or none of it does. If $B\ne0$, at most one point can vanish, with the only candidate
$I=-AB^{-1}$, which must also satisfy $I^2=-1$. This is the analytic counterpart of the
polynomial sphere remainder. It gives a local shape of the zero set, not a finiteness theorem
for the total number of zeros of every analytic function.

\subsection{A coefficientwise quaternionic Cauchy formula}\label{sec:cauchy}
Fix an \emph{ordinary} $I_0\in\mathbb S(\R)$, and a bounded ordinary conjugation-invariant
domain $U$ whose closure is contained in $\Omega$ and whose positively oriented boundary
is piecewise smooth.
Its copy in $\R+\R I_0$ is denoted $U_{I_0}$. Put
\begin{equation}\label{eq:cauchykernel}
 K(s,q)=-\bigl(q^2-2\Rea(s)q+|s|^2\bigr)^{-1}(q-\bar s).
\end{equation}
If the standard-part conjugacy sphere of $q$ comes from a point of $U$, the quadratic
denominator is nonzero for $s\in\partial U_{I_0}$. The symmetry of $U$ ensures that both
ordinary slice representatives lie in the interior rather than on its boundary.

Define an integral of a common-support Hahn family along this ordinary contour by
integrating each ordinary quaternion coefficient. This is a definition of a coefficientwise
functional; it is not a fine-topological limit of quaternion Riemann sums.

\begin{theorem}[Common-domain Cauchy formula]\label{thm:cauchy}
For a common-domain Hahn slice function as above and any finite $q$ in the standard-part
thickening of the circularization of $U$,
\begin{equation}\label{eq:cauchy}
 f(q)=\frac1{2\pi}\int_{\partial U_{I_0}}^{\!H}
 K(s,q)(-I_0\,ds)f(s).
\end{equation}
The superscript $H$ denotes the coefficientwise construction just specified.
\end{theorem}
\begin{proof}
First take ordinary $q\in U_{I_0}$. The commuting factors in its slice make
$K(s,q)=(s-q)^{-1}$. Each coefficient of $f|_{I_0}$ is an ordinary left-complex-holomorphic
quaternion-valued function; splitting its values into two complex components gives the
ordinary Cauchy formula with multiplication on the displayed sides. All ordinary derivatives
of that formula hold as well, so Taylor evaluation proves it at infinitesimally displaced
points in this fixed slice.

For a fixed general finite $q$, expand the rational kernel about $q_0=\st(q)$. On the
ordinary compact boundary its zeroth-order denominator has no zero. Its reciprocal admits
a Neumann expansion in the finitely many infinitesimal coordinates of $q-q_0$. All kernel
coefficients are ordinary smooth functions on the contour, and all supports lie in the
positive monoid generated by those infinitesimal supports. Multiplication by the common
support $S$ of $f(s)$ produces a well-ordered support with finite overlap, uniformly as a
support certificate along the contour. Thus the stated coefficientwise integral exists and
the coefficient identities may be integrated.

Finally $K(s,q)$ is a slice function of $q$: its quadratic factor has real coefficients and
its numerator is affine with the constant on the right. Apply the representation formula
to both sides of \eqref{eq:cauchy} and use the fixed-slice result at $a\pm I_0b$ to obtain
the equality at $a+Jb$ for every $J\in\mathbb S(\No)$.
\end{proof}

Only an ordinary contour in coefficient space has been used. Ordinary compactness in the
proof concerns that real contour, not the full surreal unit sphere. The kernel is the
classical slice Cauchy kernel \cite{GP}; the additional argument here supplies the Hahn
support certificate and clarifies what its integral means over surreal coefficients.

A common domain cannot be removed from the hypotheses merely by saying that every
coefficient is an analytic germ. For example,
\begin{equation}\label{eq:nocommondomain}
 F(z)=\sum_{r\ge1}\frac{t^{r\eta}}{1-rz},\qquad \eta>0,
\end{equation}
has a convergent germ for each coefficient, but its coefficient poles $1/r$ approach $0$.
There is no ordinary disk about $0$ on which all of them are holomorphic, and no common
small circle on which the preceding holomorphic-inside-contour proof applies. Formal
evaluation on a suitable infinitesimal disk can still make sense. That is a different
function class, precisely the distinction emphasized in the repository \cite{RepoDocs}.

\subsection{Fueter regularity is a different condition}\label{sec:fueter}
For quaternion coordinates $q=a+x_1\ii+x_2\jj+x_3\kk$, define the left Fueter operator and
the four-dimensional Laplacian by
\[
 \DD=\partial_a+\ii\partial_{x_1}+\jj\partial_{x_2}+\kk\partial_{x_3},
 \qquad
 \Delta_4=\partial_a^2+\partial_{x_1}^2+\partial_{x_2}^2+\partial_{x_3}^2.
\]
Coefficients multiply on the left. Here again the partial derivatives are variable
operators on ordinary coefficient functions and their Taylor extensions, not $D$.

Let $r=\sqrt{x_1^2+x_2^2+x_3^2}$ and $I=(x_1\ii+x_2\jj+x_3\kk)/r$ away from the real
axis. For an axial slice function $f=A(a,r)+I B(a,r)$ satisfying \eqref{eq:CR}, direct
coordinate differentiation gives
\begin{align}
 \DD f&=-\frac{2B}{r},\label{eq:fueterslice}\\
 \Delta_4 f&=\frac{2A_r}{r}
       +I\left(\frac{2B_r}{r}-\frac{2B}{r^2}\right).\label{eq:fuetertransform}
\end{align}
Thus a slice-regular function need not itself be Fueter-regular. The coordinate function
$f(q)=q$, for instance, has $\DD f=-2$.

\begin{theorem}[Coefficientwise Fueter mapping]\label{thm:fueter}
For a common-domain Hahn slice function constructed from holomorphic stems,
\begin{equation}\label{eq:fuetertheorem}
 \DD\Delta_4 f=0.
\end{equation}
The formula holds away from the real axis and extends across that axis wherever the stem is
defined on an ordinary neighborhood there. The extension is coefficientwise analytic.
\end{theorem}
\begin{proof}
Set $C=2A_r/r$ and $E=2B_r/r-2B/r^2$. The formula for the Fueter operator of an arbitrary
axial function is
\[
 \DD(C+IE)=(C_a-E_r-2E/r)+I(E_a+C_r).
\]
Using $A_a=B_r$, $B_a=-A_r$, and their differentiated identities gives
\begin{align*}
 C_a-E_r-2E/r
 &=\frac{2B_{rr}}r-
   \left(\frac{2B_{rr}}r-\frac{4B_r}{r^2}+\frac{4B}{r^3}\right)
   -\left(\frac{4B_r}{r^2}-\frac{4B}{r^3}\right)=0,\\
 E_a+C_r
 &=\frac{2B_{ar}}r-\frac{2B_a}{r^2}
   +\frac{2A_{rr}}r-\frac{2A_r}{r^2}=0.
\end{align*}
The parity of a stem makes $A$ even and $B$ odd in $r$. The apparent divisions in
\eqref{eq:fuetertransform} therefore have removable analytic singularities at $r=0$;
for example $B_r/r-B/r^2$ cancels the contribution of the linear term in $B$.
This proves the classical coefficient identity and its removable extension. Differentiation
of coefficient functions does not enlarge their exponent support, and Taylor evaluation is
strongly summable, so all identities extend coefficientwise to the Hahn family.
\end{proof}

This is the Fueter transformation in a specified surreal analytic setting, not a new claim
about the classical theorem \cite{GP}. Its derivation also makes the convention visible:
right-handed regularity would require the corresponding right-handed operators and product
order.

For concrete checks,
\[
 \Delta_4(q^2)=-4,\qquad
 \Delta_4(q^3)=-4\bigl(3a+x_1\ii+x_2\jj+x_3\kk\bigr),
\]
and applying $\DD$ to either right side gives zero. Multiplying such a slice polynomial by
$t^\gamma$ or taking a common-support Hahn combination preserves the coefficientwise theorem.

\subsection{What has and has not been constructed}
We now have formal infinitesimal functions, common-domain holomorphic stems, their slice
products, a coherent identity principle, a Cauchy functional, and a Fueter transformation.
The full surreal quaternion class does not thereby become an ordinary Banach-space manifold.
We have not asserted an unrestricted maximum principle, a global residue theorem on arbitrary
surreal paths, or existence of a Taylor expansion for every fine-differentiable function.
For an application requiring such a statement, its function class and integration or support
notion must be specified first.

\section{Exact computation and a formalization route}\label{sec:computation}

\subsection{A representation hierarchy}
A surquaternion is represented by four scalar objects, but the capabilities of those objects
remain decisive. No finite-string format can encode every member of the proper class $\No$.
Even within a set-sized Hahn field, an arbitrary support or coefficient sequence need not be
computable. The following hierarchy makes achievable operations explicit.

\begin{center}\small
\begin{tabular}{L{0.22\textwidth}L{0.34\textwidth}L{0.34\textwidth}}
\toprule
Scalar model & Exact quaternion operations & Additional requirements\\
\midrule
Rational functions in central monomials & Four-coordinate arithmetic, conjugation, inverse,
normal polynomials, Cayley chart & Ordered monomial group and scalar equality\\
Real algebraic extensions of that field & Modulus, square roots, finite real-closed algebra,
spectral equations & Certified algebraic roots and order\\
Effective Hahn series with support certificates & Neumann inverse, infinitesimal exponential,
logarithm, finite coefficient extraction & Decidable support queries and finite-overlap witnesses\\
Common-domain analytic coefficient families & Slice product, Taylor evaluation, coefficient
Cauchy and Fueter operators & Domain certificates and exact coefficient analytic operations\\
Unrestricted symbolic surreal expressions & Finite denotation of some further values & No
blanket decidability or evaluation guarantee\\
\bottomrule
\end{tabular}
\end{center}

The division formula $q^{-1}=\bar q/N(q)$ should normally be the first implementation path.
It needs one central inverse, not a noncommutative linear-system solver. Positive roots and
moduli should be delegated to the scalar real-closed model. A rational Hahn model is not
closed under every square root; ramification or algebraic extension may be required.

A series evaluator must return a value together with a summability certificate, or reject the
request as unsupported in that representation. It must not silently reinterpret
$\sum(\omega\ii)^n/n!$ as a global exponential. Likewise a polynomial implementation must
keep the $\star$ product separate from pointwise function multiplication. Example
\ref{ex:poly} is a compact regression test for that distinction.

Valuation error bounds are exact symbolic data. For example, truncating a Neumann inverse
after degree $M$ certifies \eqref{eq:neumannerror}; a spectral perturbation can carry both
$w(E)$ and the gap valuation $w(g)$ to certify \eqref{eq:projectorval}. These are stronger
than a floating-point residual when the important effect lives below every ordinary real
scale. They do not by themselves certify a limit in a topology whose value group has larger
scales than all multiples of the current step.

\subsection{Lean layers suggested by the repository}
The repository snapshot uses generic algebra and Hahn-series infrastructure rather than a
completed surreal construction \cite{Repo}. A compatible extension should preserve that
order of dependencies.

\paragraph{Layer A: finite quaternion algebra.}
Work with a generic ordered or real closed coefficient field, reusing the proof assistant's
quaternion algebra where suitable. Formalize the norm, inverse, centralizer, conjugacy
criterion, central-polynomial division, and normal-polynomial zero algorithm before introducing
surreal numbers. All statements quantify over a set-sized type.

\paragraph{Layer B: geometry and matrices.}
Add $F$-valued modulus, unit quaternions, the rotation action, Hermitian diagonalization,
Weyl's inequality, and the projector estimate. Avoid an ordinary real-normed-space abstraction
that would incorrectly force the modulus into $\R$. An algebraic ordered-field formulation
keeps the surreal specialization possible.

\paragraph{Layer C: quaternion Hahn series.}
Prove the four-coordinate equivalence
$\mathbb H_{\R((t^\Gamma))}\simeq\HR((t^\Gamma))$, with $\Gamma$ a type-sized ordered
abelian group. Reuse the repository's support, regrouping, and standard-part theorems.
The new obligations are ordered noncommutative coefficient convolution, leading-coefficient
noncancellation, the norm valuation identity, and the infinitesimal exponential/logarithm
substitution lemmas.

\paragraph{Layer D: the actual surreal bridge.}
Supply a separately justified surreal construction and normal-form interpretation, with an
explicit universe or class discipline. Only at this stage specialize the generic results to
surreal coefficients and prove localization. Generic Hahn lemmas alone do not establish this
bridge. No set-sized type should be advertised as containing every surreal number at the same
foundational level.

\paragraph{Layer E: analytic structures and their incompatibilities.}
Formalize common-domain stems and the coefficientwise integral using an ordinary complex
analytic library. Independently formalize the chosen surreal derivation and the properties
used in \cref{thm:nooscillation}. The no-oscillation argument is short once those substantial
foundations exist. A global trigonometric normalization belongs in another explicitly named
structure; one must not give it an unjustified instance of a chain-rule axiom.

These are proposed layers, not claims that corresponding new modules exist or compile.
The accompanying archive contains no purported Lean proof of the article. It instead
contains a standalone finite symbolic check suite, described next.

\subsection{Scope of the executable checks}
The Python program in \texttt{code/verify.py} uses exact SymPy arithmetic. It checks the
Hamilton table, generic associativity and norm identities, the complex matrix representation,
Rodrigues' formula, the derivation reconstruction formula, the polynomial example, a Cayley
unit, truncated noncommutative exponential identities, the two-scale matrix example, and the
Fueter polynomial identities. Its output is recorded in \texttt{data/verification.json}. The supplied run passed all
69 check groups, comprising 252 scalar identities, using SymPy 1.14.0.

These checks detect algebraic sign errors, misplaced factors, and expansion mistakes in the
displayed examples. They do \emph{not} verify set/class foundations, the Hahn support lemma,
real-closed-field transfer, all analytic hypotheses, or the Berarducci--Mantova construction.
The general proofs remain the mathematical arguments in this article, relying on the
explicitly cited foundational results.

\section{Further directions and limits of the present theory}\label{sec:directions}

\subsection{Collision-stable spectral subspaces}
The rank-one estimate can be developed into cluster-projector estimates with a gap between
two eigenvalue groups and no gap within either group. The natural quaternionic proof solves
an off-diagonal Sylvester equation after spectral diagonalization and estimates denominators
by the intercluster gap. Iterating this construction across a hierarchy of valuations should
produce a block-diagonal multiscale spectral calculus. The worked two-scale example proves
one rank-one instance, not that full iterative theorem or an optimal block constant.

\subsection{Finite analytic zero algebras over quaternion coefficients}
The sphere remainder algorithm completely solves finite-degree one-sided polynomials.
For common-domain analytic stems, a useful next target is a support-certified local
preparation theorem that tracks isolated and spherical zeros under infinitesimal perturbation.
Such a theorem must specify whether it takes place in a fixed-domain ring, a common-domain
germ ring, or a ring of unrelated coefficient germs. Example \eqref{eq:nocommondomain}
shows why conflating these rings can invalidate the proof before noncommutativity even enters.

\subsection{Differential extensions containing oscillations}
Theorem \ref{thm:nooscillation} gives a definite obstruction rather than a vague warning:
constant-frequency rotations are absent from the componentwise Berarducci--Mantova
quaternionic differential algebra. One concrete research direction is to adjoin solutions
of $Dy=Iy$ with a controlled involution and norm, and study what replaces the finite-element
smallness property. Another is to retain the underlying surquaternions but choose a different
derivation, making its constants and order compatibility explicit. A third is to use only
variable derivatives of common-domain stems and never identify them with a field derivation.
These are different mathematical choices, not interchangeable notations for a single calculus.

\subsection{Beyond associativity}
The Cayley--Dickson construction over $\No$ also produces an octonion algebra with positive
sum-of-eight-squares norm and therefore no zero divisors. However it is alternative rather
than associative. The centralizer, linear-factor, and matrix proofs in this article cannot
be copied without revisiting their uses of associativity. Later Cayley--Dickson algebras
already have zero divisors over $\R$ \cite{Baez}, and those examples persist under scalar extension to
$\No$. Infinite coefficients do not repair an algebraic zero-divisor identity.
This observation suggests a broader program but is not a substitute for an octonionic theory.

\subsection{The resulting picture}
Surquaternions form an associative division algebra with an ordered-field-valued norm,
algebraically closed slices, exact rotation geometry, and a robust finite spectral and
polynomial theory. Their normal forms carry a noncommutative residue field $\HR$, while the
value group records genuinely surreal scales. Positive-support calculus supplies precise
infinitesimal exponentials, logarithms, and perturbation expansions. Ordinary coefficient
analysis supplies coherent slice and Fueter operations.

The structures fit together provided their boundaries remain visible. The global phase of
an infinite imaginary argument is not determined by a raw power series. Hahn summation is
not a universal topological limit. The field derivation is not the same as differentiation of
an analytic variable. The no-oscillation theorem proves that at least one tempting attempt to
identify these structures is impossible, while the constructive results show how much useful
theory remains once that identification is abandoned.

\appendix
\section{A compact reference of conventions and dependencies}\label{app:ledger}

\begin{longtable}{L{0.23\textwidth}L{0.67\textwidth}}
\toprule
Notation or construction & Exact convention\\
\midrule
\endhead
$\HF$, $\SQ$ & Hamilton algebra over a real closed field $F$, and over $\No$, respectively.
Scalar coefficients are central.\\
$\bar q$, $N(q)$, $|q|$ & Quaternion conjugation, reduced norm $q\bar q$, and its positive
square root in the ordered center.\\
$t^\gamma$ & Conway monomial $\omega^{-\gamma}$, not an unqualified exponential power.\\
$\supp(q)$, $w(q)$ & Well-ordered set of increasing exponents; minimum exponent, with
$w(0)=+\infty$.\\
$\OO_H$, $\mm_H$, $\st$ & Finite coordinates, infinitesimal coordinates, and coordinate-zero
extraction to ordinary quaternions.\\
$P\star Q$ & Product with a central indeterminate and right-coefficient evaluation, not
pointwise multiplication.\\
$\mathbb S(F)$ & Imaginary square roots of $-1$; $F_I=F+FI$ is a complex slice.\\
$\Exp_{\mm}$, $\Log_{\mm}$ & Strong power-series operations at positive valuation.\\
$\Exp_0$ & Gonshor scalar exponential plus canonical finite circular angle, with finite
imaginary part but unrestricted scalar part.\\
$\Exp_{\mathrm{EK}}$ & Specified global extension using the initial integer part.\\
$D$ & Componentwise Berarducci--Mantova derivation, normalized by $D\omega=1$.\\
$\partial_a,\partial_b$ & Variable derivatives of ordinary coefficient functions and their
formal Taylor extensions; not $D$.\\
$\iota$ & Auxiliary central complex unit for stems; distinct from quaternion units.\\
$\DD$, $\Delta_4$ & Left Fueter operator and four-coordinate variable Laplacian.\\
$\int^H$ & Ordinary contour integration of individual Hahn coefficient functions, with a
uniform support certificate.\\
\bottomrule
\end{longtable}

\begin{center}\small
\begin{tabular}{L{0.36\textwidth}L{0.54\textwidth}}
\toprule
Results & Main dependencies beyond their displayed proofs\\
\midrule
Division algebra, centralizers, conjugacy & Ordered-field algebra; positive square roots where
normalization is used.\\
Polynomial root algorithm, spectral theorem & Real closedness of $F$ and algebraic closedness
of $F[\ii]$.\\
Surreal realization and localization & Conway normal-form theorem and real-closed Hahn-field
theorem.\\
Strong inverse, principal logarithm, BCH & Positive-support Neumann lemma and coefficientwise
strong regrouping.\\
Global quaternion exponential & Gonshor real exponential and Ehrlich--Kaplan global circular
character.\\
No-oscillation theorem & Berarducci--Mantova constant field, smallness, and $D\omega=1$ for
the chain-rule corollary.\\
Cauchy and Fueter extension & Ordinary coefficient analysis on one common domain, plus strong
Taylor support certificates.\\
\bottomrule
\end{tabular}
\end{center}

\section{Worked exact calculations for reuse}\label{app:calculations}

\subsection{Conjugating two imaginary units}
For $I=\ii$ and $J=\jj$, the conjugating element of \cref{thm:conjugacy} is
$g=1-\jj\ii=1+\kk$. Normalization gives $u=(1+\kk)/\sqrt2$, and
\[
 u\ii u^{-1}=\jj,\qquad u\jj u^{-1}=-\ii,\qquad u\kk u^{-1}=\kk.
\]
Replacing the constant direction by a surreal one does not change the formula. For example
$J=(\ii+t\jj)/\sqrt{1+t^2}$ produces an exact infinitesimal rotation through the same
construction, with no numerically approximated angle.

\subsection{An inverse and its error certificate}
For $X=t\ii$, $X^2=-t^2$, so
\[
 (1+X)^{-1}=\frac{1-t\ii}{1+t^2}
 =1-t\ii-t^2+t^3\ii+t^4-\cdots.
\]
The difference from the first $M+1$ powers $\sum_{n=0}^M(-X)^n$ is exactly
$(-X)^{M+1}(1+X)^{-1}$, of valuation $M+1$. This is an equality and a certificate, even
in the full surreal topology where the displayed partial sums do not converge as an ordinary
sequence.

\subsection{The first noncommutative correction}
For $X=t\ii$, $Y=t\jj$, the degree-two term of $\BCH(X,Y)$ is $t^2\kk$. Its degree-three
term is
\[
 \frac1{12}\bigl([t\ii,[t\ii,t\jj]]+[t\jj,[t\jj,t\ii]]\bigr)
 =-\frac{t^3}{3}(\ii+\jj).
\]
Thus
\[
 \BCH(t\ii,t\jj)
 =t(\ii+\jj)+t^2\kk-\frac{t^3}{3}(\ii+\jj)
  +\text{terms of valuation at least }4.
\]
A scalar exponential implementation that simply adds the exponents would lose the entire
$t^2\kk$ rotation. The verification program checks this coefficient calculation by truncated
noncommutative multiplication, independently of the displayed bracket simplification.

\subsection{The two-scale spectrum at integer exponents}
For $\alpha=1$, $\beta=2$, \eqref{eq:twoeigen} specializes to
\[
 \lambda_-=\frac{t-t\sqrt{1+4t^2}}2
 =-t^3+t^5-2t^7+5t^9+\cdots,
\]
\[
 \lambda_+=t+t^3-t^5+2t^7-5t^9+\cdots.
\]
The off-diagonal entry of its lower projector is
\[
 -\frac{t\ii}{\sqrt{1+4t^2}}
 =-t\ii+2t^3\ii-6t^5\ii+20t^7\ii+\cdots.
\]
The ordinary standard-part matrix is zero, yet a meaningful eigenspace is resolved at order
$t$, with its perturbation explicitly visible. This is one reason ordinary floating-point
substitution at a small positive real $t$ should be treated as an illustration rather than as
the definition of the surreal problem.

\begingroup
\small
\raggedright
\begin{thebibliography}{99}
\setlength{\itemsep}{0.35em}

\bibitem{Conway}
J. H. Conway, \emph{On Numbers and Games}, second edition, A K Peters, 2001
(first edition, Academic Press, 1976).

\bibitem{Gonshor}
H. Gonshor, \emph{An Introduction to the Theory of Surreal Numbers}, London Mathematical
Society Lecture Note Series 110, Cambridge University Press, 1986.
\href{https://doi.org/10.1017/CBO9780511629143}{doi:10.1017/CBO9780511629143}.

\bibitem{vdDE}
L. van den Dries and P. Ehrlich, Fields of surreal numbers and exponentiation,
\emph{Fundamenta Mathematicae} \textbf{167} (2001), 173--188;
erratum, \textbf{168} (2001), 295--297.

\bibitem{BM}
A. Berarducci and V. Mantova, Surreal numbers, derivations and transseries,
\emph{Journal of the European Mathematical Society} \textbf{20} (2018), 339--390.
\href{https://doi.org/10.4171/JEMS/769}{doi:10.4171/JEMS/769}.
Preprint: \href{https://arxiv.org/abs/1503.00315}{arXiv:1503.00315}.

\bibitem{EK}
P. Ehrlich and E. Kaplan, Surreal ordered exponential fields,
\emph{The Journal of Symbolic Logic} \textbf{86} (2021), no.~3, 1066--1115.
Preprint: \href{https://arxiv.org/abs/2002.07739}{arXiv:2002.07739}.

\bibitem{Voight}
J. Voight, \emph{Quaternion Algebras}, Graduate Texts in Mathematics 288, Springer, 2021.
Open-access book and corrections:
\url{https://jvoight.github.io/quat.html}.

\bibitem{GP}
R. Ghiloni and A. Perotti, Slice regular functions on real alternative algebras,
\emph{Advances in Mathematics} \textbf{226} (2011), 1662--1691.
\href{https://arxiv.org/abs/1008.4318}{arXiv:1008.4318}.

\bibitem{Baez}
J. C. Baez, The octonions, \emph{Bulletin of the American Mathematical Society}
\textbf{39} (2002), 145--205.
\href{https://arxiv.org/abs/math/0105155}{arXiv:math/0105155}.

\bibitem{Repo}
V. Reshetnikov, \emph{Surreal}, repository README and module inventory,
commit \texttt{39f2be6667ade51bca2b45daa47e289d69c09764}, accessed September 21, 2026.
\url{https://github.com/VladimirReshetnikov/Surreal}.

\bibitem{RepoDocs}
V. Reshetnikov, \emph{The surreal and surcomplex reports},
\texttt{docs/README.md} in the same repository and commit.
\href{https://github.com/VladimirReshetnikov/Surreal/blob/39f2be6667ade51bca2b45daa47e289d69c09764/docs/README.md}{Pinned catalogue source}.

\bibitem{RepoTrig}
\emph{Trigonometry on the Surcomplex Plane: Canonical Phases, Arbitrary-Scale Geometry,
and Infinitesimal Degeneration}, unpublished merged report in the same repository,
\texttt{docs/surcomplex/trigonometry/article.tex}, September 21, 2026.
\href{https://github.com/VladimirReshetnikov/Surreal/blob/39f2be6667ade51bca2b45daa47e289d69c09764/docs/surcomplex/trigonometry/article.tex}{Pinned trigonometry source}.

\end{thebibliography}
\endgroup
\end{document}
